
\chapter{Applications of the trace formula to trigonometric sums}\label{VI}
\chaptermark{Applications to trigonometric sums}










In this exposé, I explain how the trace formula can be used to calculate 
or study various trigonometric sums and how, together with the 
Weil conjecture, it can be used to bound them.

The first two sections provide a ``user's guide'' to these tools. 
\hyperref[VI:3]{Section 3} presents, in cohomological language, Weil's results on sums in 
$1$ variable. Sections $4$ through $6$ give a 
detailed study of Gauss and Jacobi sums -- including 
Weil's old and recent results on Hecke characters 
defined by Jacobi sums. In Section $7$, we study a 
generalization of Kloosterman sums to several variables. Finally, 
Section $8$ gives some indications of other uses that have been, or can be, 
made of these methods. 










\section*{Notations}\label{VI:0}





\subsection{}\label{VI:0-1}

We use the notation of \hyperref[II]{Report}, Section \ref{II:1}. We 
will often have to consider a finite extension $\dF_{q^n}\subset \dF$ of 
$\dF_q$. A subscript $0$ will denote an object over $\dF_q$, and a subscript $1$ 
an object over $\dF_{q^n}$. Replacing a subscript $0$ by a subscript $1$ (resp. 
omitting the subscript) means extending scalars to $\dF_{q^n}$ (resp. 
to $\dF$).





\subsection{}\label{VI:0-2}

We denote by $\ell$ a prime number $\ne p$. We freely use the 
language of $\dQ_\ell$-sheaves, as well as that of $E_\lambda$-sheaves, 
where $E_\lambda$ is a finite extension of $\dQ_\ell$ (cf. \hyperref[II]{Report}, 
Section \ref{II:2}, especially \ref{II:2-11}). 





\subsection{}\label{VI:0-3}

Let $H^\bullet$ be a graded vector space. If $T$ is an endomorphism of 
$H^\bullet$, we put (cf. \hyperref[IV]{Cycle}, \ref{IV:eq:1-3-5}) 
\[
  \tr(T,H^\bullet) = \sum (-1)^i \tr(T,H^i) \text{.}
\]










\section{Principles}\label{VI:1}





\subsection{}\label{VI:1-1}

Let $X_0$ be a separated scheme of finite type over $\dF_q$, let $E_\lambda$ be a 
finite extension of $\dQ_\ell$, and let $\sF_0$ be an $E_\lambda$-sheaf on $X_0$. The 
trace formula says that 
\begin{equation*}\tag{1.1.1}\label{VI:eq:1-1-1}
  \sum_{x\in X^F} \tr(F_x^\ast,\sF) = \sum (-1)^i \tr\left(F^\ast,\h_c^i(X,\sF)\right) \text{.}
\end{equation*}

With the notation \ref{VI:0-3}, the right-hand side is simply written 
$\tr(F^\ast,\h_c^\bullet(X,\sF))$. 

This trace formula for $E_\lambda$-sheaves can either be 
deduced from \hyperref[II]{Report} \ref{II:4-10} by passage to the limit (cf. 
\hyperref[II]{Report} \ref{II:4-11} through \ref{II:4-13}), or be deduced 
from the trace formula for $\dQ_\ell$-sheaves (\hyperref[II]{Report} 
\ref{II:3-2}) by the method of (\hyperref[III]{Functions $L$ mod. $\ell^n$}, 
\ref{III:4-3}). 

We shall interpret various trigonometric sums as the left-hand side of 
\eqref{VI:eq:1-1-1}, for suitable $X_0$ and $\sF_0$. 





\subsection{}\label{VI:1-2}

Let $A$ be a finite group. An \emph{$A$-torsor} on a scheme $X$ is a 
sheaf $T$ on $X$, equipped with a right action of $A$, which, locally 
(for the étale topology) on $X$, is isomorphic to the constant sheaf $A$ on 
which $A$ acts by right translation. 

If $\tau:A\to B$ is a homomorphism and $T$ is an $A$-torsor, there exists, 
up to unique isomorphism, a unique $B$-torsor $\tau(T)$, equipped with a 
morphism of sheaves $\tau:T\to \tau(T)$ such that 
\begin{equation*}\tag{1.2.1}\label{VI:eq:1-2-1}
  \tau(t a) = \tau(t) \tau(a)\text{.}
\end{equation*}

If $T$ is an $A$-torsor on $X$ and $\rho$ is a linear representation 
of $A$, $\rho:A\to \gl(V)$ (where $V$ is a finite-dimensional vector space over 
$E_\lambda$), there exists, up to unique isomorphism, a unique 
$E_\lambda$-sheaf $\sF$, smooth of rank $\dim(V)$, equipped with a morphism of 
sheaves 
\begin{equation*}\tag{1.2.2}\label{VI:eq:1-2-2}
  \rho:T\to \underline{\operatorname{Isom}}(V,\sF) \text{.}
\end{equation*}
such that $\rho(t a) = \rho(t) \rho(a)$. We denote it by $\rho(T)$. For 
$\tau:A\to B$ and $\rho$ a representation of $B$, we have a canonical 
isomorphism $\rho(\tau(T)) = (\rho\tau)(T)$. 

In this paragraph (except the appendix), we only consider the case where $A$ 
is commutative and where $V=E_\lambda$: $\rho$ is a character 
$A\to E_\lambda^\times$, and $\rho(T)$ is a smooth $E_\lambda$-sheaf of rank 
one. The morphism \eqref{VI:eq:1-2-2} is interpreted as an everywhere 
nonzero morphism 
\begin{equation*}\tag{1.2.3}\label{VI:eq:1-2-3}
  \rho:T\to \rho(T) 
\end{equation*}
such that $\rho(t a) = \rho(a)\rho(t)$. 





\subsection{}\label{VI:1-3}

Let $S$ be a scheme, let $G$ be a commutative group scheme over $S$, and let $G'$ 
be an extension of commutative group schemes over $S$
\[\xymatrix{
  0 \ar[r] 
    & A \ar[r] 
    & G' \ar[r]^-\pi 
    & G \ar[r] 
    & 0 \text{.}
}\]
The sheaf $T$ of local sections of $\pi$ is then an $A$-torsor on 
$G$. 

If $X$ is a scheme over $S$ and $f,g$ are two $S$-morphisms from $X$ to $G$, 
the fact that $T$ is defined by an extension gives a canonical isomorphism 
\begin{equation*}\tag{1.3.1}\label{VI:eq:1-3-1}
  (f+g)^\ast T = f^\ast T + g^\ast T \text{.}
\end{equation*}
On the left, $f+g$ is the sum of $f$ and $g$ with respect to the group law on $G$; 
on the right, $+$ denotes a sum of torsors. 

If $f:X\to G$ factors through $G'$, $f^\ast T$ is trivial (and a 
factorization gives a trivialization); up to (non-unique) isomorphism, 
the $A$-torsor $f^\ast T$ depends only on the image of $f$ in 
$\hom_S(X,G)/\pi\hom_S(X,G')$: it is given by the morphism $\partial$ in 
the exact sequence 
\[\xymatrix{
  \hom(X,G') \ar[r] 
    & \hom(X,G) \ar[r]^-\partial 
    & \h^1(X,A) \text{.}
}\]

We will apply these constructions to the exact Kummer sequences 
$0 \to \dmu_n \to \dG_m \to \dG_m \to 0$, and to Lang torsors. 





\subsection{The Lang torsor}\label{VI:1-4}

Let $G_0$ be a connected commutative algebraic group over $\dF_q$ (for the 
noncommutative case, see the appendix). The group law is written 
\emph{multiplicatively}. Lang's isogeny 
\[
\fL:G_0 \to G_0: x\mapsto F x\cdot x^{-1}
\]
is étale; its image, an open subgroup of $G_0$, can only be 
$G_0$ itself; its kernel is the finite group $G_0(\dF_q)$. The \emph{Lang 
torsor} $L$ is the $G_0(\dF_q)$-torsor on $G_0$ defined by the exact sequence 
\begin{equation*}\tag{1.4.1}\label{VI:eq:1-4-1}
\xymatrix{
  0 \ar[r] 
    & G_0(\dF_q) \ar[r] 
    & G_0 \ar[r]^-\fL 
    & G_0 \ar[r] 
    & 0 \text{.}
}
\end{equation*}

\paragraph{Notation}
We also write $L$ and $\fL$ (or $L_{(q)}$ and $\fL_{(q)}$) for the objects obtained 
from $L$ and $\fL$ by extending the base field. When necessary, we 
specify them with a subscript $0$ or $1$.


\paragraph{Examples}
For $G_0=\dG_a$ or $\dG_m$, the sequences \eqref{VI:eq:1-4-1} are written 
\begin{align*}\tag{1.4.2}\label{VI:eq:1-4-2}
&\xymatrix{
  0 \ar[r] 
    & \dF_q \ar[r] 
    & \dG_a \ar[r]^-{x^q-x} 
    & \dG_a \ar[r] 
    & 0 
} \\
\tag{1.4.3}\label{VI:eq:1-4-3}
&\xymatrix{
  0 \ar[r] 
    & \dmu_{q-1} \ar[r] 
    & \dG_m \ar[r]^-{x^{q-1}} 
    & \dG_m \ar[r] 
    & 0
}
\end{align*}





\subsection{}\label{VI:1-5}

Let us compute the endomorphism $F^\ast$ of the fiber of the Lang torsor at 
$\gamma\in G^F=G_0(\dF_q)$. If $g\in \fL^{-1}(\gamma)\subset G$, we have 
$F g = F g\cdot g^{-1}\cdot g = \gamma g$. Therefore (\hyperref[II]{Report}, 
\ref{II:1-2}) 
\begin{equation*}\tag{1.5.1}\label{VI:eq:1-5-1}
  \text{on $L_0(G_0)_\gamma\simeq \fL^{-1}(\gamma)$, $F^\ast$ is $g\mapsto g\gamma^{-1}$.}
\end{equation*}





\subsection{}\label{VI:1-6}

On $\dF_{q^n}$, the identity $F_{(q^n)}=F_{(q)}^n$ between endomorphisms of 
$G_1$ implies that $\fL_{(q^n)}=\fL_{(q)}\circ \prod_{i=0}^{n-1} F_{(q)}^i$. 
On $G_0(\dF_{q^n})$, $F_{(q)}^i$ acts as the element 
$x\mapsto x^{q^i}$ of $\gal(\dF_{q^n}/\dF_q)$. On $G_0(\dF_{q^n})$, 
$\prod F_{(q)}^i$ is therefore the composite of the norm 
$G_0(\dF_{q^n})\to G_0(\dF_q)$ and the inclusion 
$G_0(\dF_q)\subset G_0(\dF_{q^n})$: the diagram 
\begin{equation*}\tag{1.6.1}\label{VI:eq:1-6-1}
\xymatrix{
  0 \ar[r] 
    & G_0(\dF_{q^n}) \ar[r] \ar[d]^-N 
    & G_1 \ar[r]^-{\fL_{(q^n)}} \ar[d]^-{\prod_{i=0}^{n-1} F_{(q)}^i} 
    & G_1 \ar[r] \ar@{=}[d] 
    & 0 \\
  0 \ar[r] 
    & G_0(\dF_q) \ar[r] 
    & G_1 \ar[r]^-{\fL_{(q)}} 
    & G_1 \ar[r] 
    & 0
}
\end{equation*}
is commutative. In the language of torsors, this gives a canonical isomorphism 
between $G_0(\dF_q)$-torsors on $G_1$ 
\begin{equation*}\tag{1.6.2}\label{VI:eq:1-6-2}
  N L_{(q^n)} = L_{(q)} \text{.}
\end{equation*}





\begin{definition_}\label{VI:1-7}
Let $f_0:X_0\to G_0$ be a morphism and 
$\chi:G_0(\dF_q)\to E_\lambda^\times$ a character. We set 
$\sF(\chi,f_0) = \chi^{-1}(f_0^\ast (L_0(G_0))) = f_0^\ast\chi^{-1}(L_0(G_0))$. 
\end{definition_}

We have the following functoriality properties:
\begin{enumerate}[a)]
  \item $\sF(\chi,f_0)$ is bimultiplicative in $\chi$ and $f_0$: 
    \begin{align*}\tag{1.7.1}\label{VI:eq:1-7-1}
      \sF(\chi,f_0'\cdot f_0'') &= \sF(\chi,f_0')\otimes \sF(\chi,f_0'') \text{,} \\
      \tag{1.7.2}\label{VI:eq:1-7-2}
      \sF(\chi'\chi'',f_0) &= \sF(\chi',f_0)\otimes \sF(\chi'',f_0) \text{.}
    \end{align*}
    This follows from the definition of $\chi(T)$, for $T$ a torsor, combined 
    with \eqref{VI:eq:1-3-1} for \eqref{VI:eq:1-7-1}. 
  \item For $g_0:Y_0\to X_0$ a morphism, we have 
    \begin{equation*}\tag{1.7.3}\label{VI:eq:1-7-3}
      \sF(\chi,f_0\circ g_0) = g_0^\ast \sF(\chi,f_0) \text{.}
    \end{equation*}
    In particular, $\sF(\chi,f_0) = f_0^\ast \sF(\chi)$, where $\sF(\chi)$ 
    denotes $\sF(\chi,\operatorname{id}_{G_0})$. 
  \item For $u_0:G_0\to H_0$ a morphism, and 
    $\chi:H_0(\dF_q) \to E_\lambda^\times$, we have 
    \begin{equation*}\tag{1.7.4}\label{VI:eq:1-7-4}
      \sF(\chi,u_0 f_0) = \sF(\chi u_0,f_0) \text{.}
    \end{equation*}
  \item For $G_0=\prod_{i\in I} G_0^i$, with $\chi$ having coordinates $\chi_i$ 
    ($i\in I$) and $f_0$ having coordinates $f_0^i$, it follows from 
    \eqref{VI:eq:1-7-1} through \eqref{VI:eq:1-7-4} that 
    \begin{equation*}\tag{1.7.5}\label{VI:eq:1-7-5}
      \sF(\chi,f_0) = \bigotimes_{i\in I} \sF(\chi_i,f_0^i) \text{.}
    \end{equation*}    
\end{enumerate}

Note the $\chi^{-1}$ in Definition \ref{VI:1-7}. It ensures that, for 
$x\in X^F$, on $\sF(\chi,f_0)_x$ one has 
\begin{equation*}\tag{1.7.6}\label{VI:eq:1-7-6}
  F_x^\ast = \chi f_0(x)
\end{equation*}
(use the fact that the fiber at $x$ of the morphism \eqref{VI:eq:1-2-2}, for 
$\rho=\chi^{-1}$, commutes with $F_x^\ast$, together with \eqref{VI:eq:1-5-1}). 

If $\sF(\chi,f_0)_1$ is the sheaf obtained from $\sF(\chi,f_0)$ by extension 
from the base field to $\dF_{q^n}$, we deduce from \eqref{VI:eq:1-6-2} that 
\begin{equation*}\tag{1.7.7}\label{VI:eq:1-7-7}
  \sF(\chi,f_0)_1 = \sF(\chi\circ N,f_1) \text{.}
\end{equation*}





\subsection{Abuse of notations}\label{VI:1-8}

(i) If $\Xi$ denotes the composite map 
$\chi f_0:X_0(\dF_q) \to E_\lambda^\times$, we will sometimes write $\sF(\Xi)$ 
instead of $\sF(\chi,f_0)$. By \eqref{VI:eq:1-7-1} through \eqref{VI:eq:1-7-5}, 
there is little risk of ambiguity. For example:
\begin{enumerate}[a)]
  \item we write $\sF(\chi)$ for $\sF(\chi,\operatorname{id}_{G_0})$ 
    (notation already used in \ref{VI:1-7}b)); 
  \item we write $\sF(\chi f_0)$ for $\sF(\chi,f_0)$;
  \item with the notation of \ref{VI:1-7}d), we write 
    $\sF(\prod \chi_i f_0^i)$ for $\sF(\chi,f_0)$. 
\end{enumerate}

With this notation, \eqref{VI:eq:1-7-4} shows that the notation 
$\sF(\chi u_0 f_0)$ is unambiguous, \eqref{VI:eq:1-7-1} \eqref{VI:eq:1-7-2} 
\eqref{VI:eq:1-7-5} express the multiplicativity of $\sF(\Xi)$ with respect to $\Xi$, 
and \eqref{VI:eq:1-7-6} becomes $F_x^\ast = \Xi(x)$ on $\sF(\Xi)_x$.

(ii) If $X$ is a scheme over an extension $k$ of $\dF_q$, and $f$ is a 
morphism of $X$ into $G_0\otimes_{\dF_q} k$, we shall again denote by $\sF(\chi,f)$, 
$\sF(\chi f)$, or $\sF(\chi)$ (for $f$ an inclusion) the inverse image under 
$f$ of the $E_\lambda$-sheaf obtained from $\chi^{-1}(L_0(G))$ by extension of 
scalars from $\dF_q$ to $k$. If $f_0$ is the composite 
$X\to G_0\otimes_{\dF_q} k \to G_0$, this is again the $\sF(\chi,f_0)$ of 
\ref{VI:1-7}. If $k$ is a finite extension of $\dF_q$, we have 
$\sF(\chi f) = \sF(\chi N_{k/\dF_q} f)$. 

Applying \eqref{VI:eq:1-1-1} to \eqref{VI:eq:1-7-6} and \eqref{VI:eq:1-7-7}, 
we find: 





\begin{theorem_}\label{VI:1-9}
Let $S_0$ be a separated scheme of finite type over $\dF_q$, let $G_0$ be a 
connected commutative algebraic group over $\dF_q$, let $f_0:S_0\to G_0$ be a morphism, and let 
$\chi:G_0(\dF_q) \to E_\lambda^\times$. We have 
\begin{equation*}\tag{1.9.1}\label{VI:eq:1-9-1}
  \sum_{s\in S_0(\dF_q)} \chi f_0(s) = \tr(F^\ast,\h_c^\bullet(S,\sF(\chi,f_0))) 
\end{equation*}
and, for any integer $n\geqslant 1$,
\begin{equation*}\tag{1.9.2}\label{VI:eq:1-9-2}
  \sum_{s\in S_0(\dF_{q^n})} \chi\circ N_{\dF_{q^n}/\dF_q} f_0(s) = \tr({F^\ast}^n, \h_c^\bullet\left(S,\sF(\chi,f_0))\right) \text{.}
\end{equation*}
\end{theorem_}





\addtocounter{subsubsection}{2}
\subsubsection{Remark}\label{VI:1-9-3}

Take $G_0$ to be a product. With the notation of \ref{VI:1-7}d) and 
\ref{VI:1-8}c), the formula \eqref{VI:eq:1-9-2} becomes 
\[
  \sum_{s\in S_0(\dF_{q^n})} \prod_i \chi_i\circ N_{\dF_{q^n}/\dF_q} (f_0^i(s)) = \tr\Big({F^\ast}^n, \h_c^\bullet\Big(S,\sF\Big(\prod_i \chi_i f_0^i\Big)\Big)\Big) \text{.}
\]





\subsubsection{Remark}\label{VI:1-9-4}

Take $G_0=\{e\}$. Formula \eqref{VI:eq:1-9-1} becomes 
\[
  \left|S_0(\dF_q)\right| = \sum_{s\in S_0(\dF_q)} 1 = \tr(F^\ast,\h_c^\bullet(S,\dQ_\ell)) \text{.}
\]

It is rare to be able to compute the right-hand side of 
\eqref{VI:eq:1-9-1} explicitly. Here is an amusing example, with $G_0=\{e\}$. 





\subsection{Example}\label{VI:1-10}

A nonsingular projective quadric of odd dimension over \texorpdfstring{$\dF_q$}{Fq} has the same number of rational points as projective space of the same dimension over \texorpdfstring{$\dF_q$}{Fq}.

If, over $\dC$, $X$ is a nonsingular quadric hypersurface in 
the projective space $\dP^{2 N}$, and $Y$ a hyperplane, we know that for 
$i\leqslant 2\dim(X)=2\dim(Y)$, the inclusions 
$X\hookrightarrow \dP^{2 N}\hookleftarrow Y$ induce isomorphisms (in 
ordinary cohomology) 
\[\xymatrix{
  \h^i(X,\dQ) 
    & \ar[l]_-\sim \h^i(\dP^{2 N},\dQ) \ar[r]^-\sim 
    & \h^i(Y,\dQ) \text{.}
}\]

By specialization, it follows that if $X_0'$ is a nonsingular quadric 
hypersurface in projective space $\dP_0^{2 N}$ over $\dF_q$, 
and $Y_0'$ a hyperplane, the inclusions 
$X_0'\hookrightarrow \dP_0^{2 N}\hookleftarrow Y_0'$ induce 
isomorphisms 
\[\xymatrix{
  \h^i(X',\dQ_\ell) 
    & \ar[l]_-\sim \h^i(\dP^{2 N},\dQ_\ell) \ar[r]^-\sim 
    & \h^i(Y',\dQ_\ell) \text{.}
}\]
These isomorphisms commute with $F^\ast$, and the trace formula applies. 





\subsection{}\label{VI:1-11}

The trace formula can also be used to understand the classical formulas 
giving the number of rational points of linear groups over 
finite fields, and of certain homogeneous spaces (see \hyperref[VI:8]{Section 8}). 





\subsection{}\label{VI:1-12}

There is a dictionary that translates various kinds of classical 
manipulations of trigonometric sums into cohomological terms. This dictionary 
will be given in \hyperref[VI:2]{Section 2}. Since the cohomological statement is more 
``geometric,'' it can sometimes be applied 
to situations where the classical argument only applies after a 
extension of the finite base field. Here is an example (a special case of a 
theorem of Kazhdan).

\begin{theorem_}[Kazhdan]\label{VI:1-13}
Let $X_0$ be a scheme over $\dF_q$. Suppose there is an action $\rho$ of 
$\dG_a$ on $X$, and a morphism $f:X\to Y$ of schemes over $\dF$, making 
$X$ a $\dG_a$-torsor over $Y$. Then the number of rational points of $X_0$ 
is divisible by $q$.
\end{theorem_}

First suppose that $\rho$, $Y$, and $f$ are defined over $\dF_q$. 


\paragraph{Classical argument}
The fibers of $f:X_0(\dF_q)\to Y_0(\dF_q)$ all have $q$ elements: 
$|X_0(\dF_q)| = q\cdot |Y_0(\dF_q)|$, hence the divisibility. 


\paragraph{Cohomological translation}
The higher direct image sheaves $\eR^i f_!\dQ_\ell$ are 
\[
  \eR^i f_!\dQ_\ell = \begin{cases}
                       \dQ_\ell(-1) & \text{if $i = 2$} \\
                       0            & \text{for $i\ne 2$}
                     \end{cases}
\]
The Leray spectral sequence for $f$ therefore degenerates, giving an isomorphism 
\[
  \h_c^i(X,\dQ_\ell) = \h_c^{i-2}(Y,\dQ_\ell)(-1) \text{.}
\]
To understand the effect of a Tate twist on the eigenvalues of 
Frobenius, it is most convenient to use the Galois point of view 
(\hyperref[II]{Report} \ref{II:1-8}) and note that the Frobenius substitution 
$\varphi$ acts on $\dQ_\ell(1)$ by multiplication by $q$. The 
eigenvalues of $F^\ast$ on $\h_c^p(X,\dQ_\ell)$ are therefore obtained by 
multiplying by $q$ the eigenvalues of $F^\ast$ acting on $\h_c^{i-2}(Y,\dQ_\ell)$. We 
know that the latter are algebraic integers \cite[XXI 5.2.2]{sga7}; hence 

\begin{equation*}\tag{$*$}\label{VI:eq:*}
  \substack{\displaystyle\text{the eigenvalues of $F^\ast$ acting on $\h_c^i(X,\dQ_\ell)$ are} \\
  \displaystyle\text{algebraic integers divisible by $q$.}}
\end{equation*}


\paragraph{Descent}
Return to the hypotheses of the theorem. For a suitable $n$, we may 
suppose that $\rho$, $Y$, and $f$ are defined over $\dF_{q^n}$. Since the 
Frobenius morphism relative to $\dF_{q^n}$ is the $n$-th power of the one 
relative to $\dF_q$, statement \eqref{VI:eq:*} for the scheme over 
$\dF_{q^n}$ obtained from $X_0$ by extension of scalars gives: 
\begin{equation*}\tag{$**$}\label{VI:eq:**}
  \substack{\displaystyle\text{the $n$-th powers of the eigenvalues of $F^\ast$ acting on} \\ \displaystyle\text{$\h_c^i(X,\dQ_\ell)$ are algebraic integers divisible by $q^n$.}}
\end{equation*}

This assertion implies \eqref{VI:eq:*}. The number of rational points of 
$X_0$, namely $\sum (-1)^i \tr(F^\ast,\h^i(X,\dQ_\ell))$, is therefore an 
algebraic integer divisible by $q$. It follows that, as a rational integer, 
it is divisible by $q$. 





\subsection{}\label{VI:1-14}

Formula \ref{VI:1-9-3} controls the dependence on $n$ of the trigonometric 
sum on the left-hand side. It implies identities between a 
trigonometric sum and those which are deduced from it by ``extension of 
scalars.'' The most famous of these identities is that of 
Hasse--Davenport. Let $\chi:\dF_q^\times \to \dC^\times$ and 
$\psi:\dF_q\to \dC^\times$ be characters, with $\psi$ nontrivial, and 
define the Gauss sum $\tau(\chi,\psi)$ by 
\begin{equation*}\tag{1.14.1}\label{VI:eq:1-14-1}
  \tau(\chi,\psi) = - \sum_{x\in \dF_q^\times} \psi(x) \chi^{-1}(x) \text{.}
\end{equation*}





\begin{theorem_}[Hasse-Davenport]\label{VI:1-15}
We have 
\begin{equation*}\tag{1.15.1}\label{VI:eq:1-15-1}
  \tau\left(\chi\circ N_{\dF_{q^n}/\dF_q},\psi\circ \tr_{\dF_{q^n}/\dF_q}\right) = \tau(\chi,\psi)^n \text{.}
\end{equation*}
\end{theorem_}

Let $E\subset \dC$ be a number field containing the values of $\chi$ and 
$\psi$, and let $E_\lambda$ be the completion of $E$ at a place of residue 
characteristic $\ell\ne p$. Identity \eqref{VI:eq:1-15-1} is an identity in 
$E$. Proving it in $\dC$ is equivalent to proving it in $E_\lambda$. We 
shall prove it in $E_\lambda$, regarding $\chi$ and $\psi$ as taking values 
in $E_\lambda^\times$.

Apply \ref{VI:1-9-3} with $X_0=\dG_m$ over $\dF_q$ and 
$G_0=\dG_m\times \dG_a$. The groups $\h_c^i(\dG_m,\sF(\chi^{-1}\psi))$ vanish for 
$i\ne 1$, and $\h_c^1$ has dimension $1$ (\ref{VI:4-2}). By 
\ref{VI:1-9-3}, 
$\tau(\chi\circ N_{\dF_{q^n}/\dF_q},\psi\circ \tr_{\dF_{q^n}/\dF_q})$ is 
the unique eigenvalue of $(F^\ast)^n$ on $\h_c^1$, and \eqref{VI:eq:1-15-1} 
follows. 

Similar examples are given in Weil \cite[App V]{we74}. 





\subsection{}\label{VI:1-16}

An $E_\lambda$-sheaf $\sF_0$ on $X_0$ is said to be \emph{pointwise pure of 
weight $n$} if for any closed point $x\in |X_0|$, the eigenvalues of 
$F_x^\ast$ (\hyperref[II]{Report} \ref{II:1-2}) are algebraic numbers 
all of whose complex conjugates have absolute value 
$q_x^{n/2}$, where $q_x$ is the number of elements of $k(x)$. The 
following theorem will be proved in \cite{de80}. 





\begin{theorem_}\label{VI:1-17}
If $\sF_0$ is pointwise pure of weight $n$, then for any eigenvalue $\alpha$ 
of the endomorphism $F^\ast$ of $\h_c^i(X,\sF)$, there exists an integer 
$m\leqslant n+i$ such that the complex conjugates of $\alpha$ are all of 
absolute value $q^{m/2}$.
\end{theorem_}

The sheaves considered in \ref{VI:1-9} have weight $0$. The theorem 
therefore gives, for the sum \ref{VI:eq:1-9-1} (or rather, for all its 
complex conjugates), the bound 
\begin{equation*}\tag{1.17.1}\label{VI:eq:1-17-1}
  \left|\sum_{s\in S_0(\dF_q)} \prod_i \chi_i(f^i(s)) \right| \leqslant \sum_i \dim \h_c^i\left(S,\sF\left(\prod \chi_i f^i\right)\right)\cdot q^{i/2} \text{.}
\end{equation*}





\subsection{Remarks}\label{VI:1-18}

\begin{enumerate}[a)]
  \item If $\dim S=n$, for any $E_\lambda$-sheaf $\sF$ on $S$, we have 
    \[
      \h_c^i(S,\sF) = 0 \qquad \text{for $i\notin [0,2n]$.}
    \]
  \item The group $\h_c^0(X,\sF)$ is the group of global sections of $\sF$ 
    on $S$ whose support is proper. If $S$ is connected, not proper and 
    $\sF$ is smooth, this group is zero; the same holds if $S$ is connected, 
    $\sF$ is smooth of rank one, and is nonconstant.
  \item If $S$ is smooth and purely of dimension $n$, and $\sF$ is smooth, 
    Poincaré duality says that the vector spaces 
    $\h_c^i(X,\sF)$ and $\h^{2n-i}(X,\sF^\vee)(n)$ are dual to each other 
    (for the effect of a Tate twist on the eigenvalues of Frobenius, 
    cf. the $2$nd part of the proof of \ref{VI:1-13}).
  \item If $\dim S=n$, Poincaré duality allows us to compute 
    $\h_c^{2n}(S,\sF)$ as follows: take an open subset $U$ of $S_\text{red}$ 
    whose complement has dimension $<n$; require this open subset to be smooth and 
    purely of dimension $n$, with $\sF$ smooth on it. We then have 
    $\h_c^{2n}(U,\sF) \iso \h_c^{2n}(S,\sF)$, because in the long exact sequence 
    of cohomology, $\h_c^i(S\setminus U,\sF) = 0$ for $i=2n,2n-1$. 
    Poincaré duality therefore gives 
    \[
      \h_c^{2n}(S,\sF) \iso \left(\h^0(U,\sF^\vee)(n)\right)^\vee \text{.}
    \]
    
    Let $U$ be connected, and let $u$ be a geometric point of $U$. The 
    ``local system'' $\sF$ corresponds to a representation of 
    $\pi_1(U,u)$ on $\sF_u$, and $\h^0(U,\sF^\vee)$ is 
    $(\sF_u^\vee)^{\pi_1(U,u)}=((\sF_u)_{\pi_1(U,u)})^\vee$ (dual of 
    coinvariants). We therefore have 
    \[
      \h_c^{2n}(S,\sF) \simeq (\sF_u)_{\pi_1(U,u)}(-n) \text{.}
    \]
  \item These remarks often allow the calculation of $b_0$ and $b_{2n}$. 
    Computing the other $b_i$ can be difficult. If $b_{2n}=0$ and 
    $b_{2n-1}=0$, the bound \eqref{VI:eq:1-17-1} is a Lang--Weil bound, 
    with, for large $q$, a gain of $\sqrt q$ over the 
    trivial $O(q^n)$ bound. One should bear in mind, however, that 
    \eqref{VI:eq:1-17-1} does not formally imply the trivial bound 
    \[
      \left|\sum_{s\in S_0(\dF_q)} \prod_i \chi_i(f^i(s))\right| \leqslant \left| S_0(\dF_q)\right| \text{,}
    \]
    so it may be useful to take this into account (cf. the proof of \ref{VI:3-8}). 
\end{enumerate}

Below I explain a method for proving that $b_i=0$ for $i\ne n$. 
When it applies, it gives an estimate of order $O(q^{n/2})$. The constant 
implicit in $O$ is $b_n=(-1)^n \chi(S,\sF(\prod \chi_i f^i))$, which may 
be difficult to compute. 





\begin{proposition_}\label{VI:1-19}
Let $\bar X$ be a proper scheme over an algebraically closed field $k$, 
let $j:X\hookrightarrow \bar X$ be an open immersion with domain $X$, and let $\sF$ be an $E_\lambda$-sheaf 
on $X$. Assume that 
\begin{enumerate}[\indent a)]
  \item $(j_\ast\sF)_x = 0$ for $x\in \bar X\setminus X$, i.e. 
    $j_!\sF\iso j_\ast\sF$; 
  \item $\eR^i j_\ast \sF = 0$ for $i>0$.
\end{enumerate}
Then the maps 
\[
  \h_c^i(X,\sF) \to \h^i(X,\sF)
\]
are isomorphisms. 
\end{proposition_}

Assumptions a) and b) mean that $j_! \sF \iso \eR j_\ast \sF$, hence 
\[
  \h_c^i(X,\sF) = \h^i(\bar X,j_!\sF) \iso \hh^i(\bar X,\eR j_\ast \sF) = \h^i(X,\sF) \text{.}
\]
In other words, the Leray spectral sequence for $j$ 
\[
  E_2^{pq} = \h^p(\bar X,\eR^q j_\ast \sF) \Rightarrow \h^{p+q}(X,\sF) 
\]
degenerates, by b) into isomorphisms 
\[
  \h^i(\bar X,j_\ast\sF) \iso \h^i(X,\sF) \text{,}
\]
while, by a) 
\[
  \h_c^i(X,\sF) = \h^i(\bar X,j_!\sF) \iso \h^i(\bar X,j_\ast\sF) \text{.}
\]





\subsubsection{Example}\label{VI:1-19-1}

If $X$ is a curve, or if $\bar X$ is smooth, $X$ is the complement of a 
normal-crossings divisor, and $\sF$ is tamely ramified, 
hypothesis a) of \ref{VI:1-19} (no invariants under local monodromy) 
implies hypothesis b). 





\begin{proposition_}\label{VI:1-20}
Let $X_0$ be a smooth affine scheme, purely of dimension $n$, over $\dF_q$, and let 
$\sF_0$ be a smooth $E_\lambda$-sheaf, pointwise pure of weight $m$, on $X_0$. We 
assume that the maps $\h_c^i(X,\sF) \to \h^i(X,\sF)$ are 
isomorphisms (cf. \ref{VI:1-17}). Then, $\h_c^i(X,\sF) = 0$ for $i\ne n$, and 
the complex conjugates of the eigenvalues of $F^\ast$ on $\h_c^n(X,\sF)$ 
are of absolute value $q^{(m+n)/2}$. 
\end{proposition_}

Poincaré duality pairs $\h_c^i(X,\sF)$ (resp. 
$\h_c^i(X,\sF^\vee)$) with $\h^{2n-i}(X,\sF^\vee(n))$ (resp. 
$\h^{2n-i}(X,\sF(n))$). According to \cite[XIV 3.2]{sga4}, the ordinary $\h^i$ groups 
vanish for $i>n$. By duality, the $\h_c^i$ are zero for $i<n$; since 
$\h_c^i(X,\sF) \iso \h^i(X,\sF)$, these groups are zero for $i\ne n$. 
Apply \ref{VI:1-17} to $\sF_0$ and to $\sF_0^\vee(n)$ (of weight 
$-m-2n$). We find that the complex conjugates $\alpha$ of the eigenvalues 
of $F^\ast$ on $\h_c^n(X,\sF) = \h_c^n(X,\sF^\vee(n))^\vee$ satisfy 
\begin{align*}
  |\alpha| &\leqslant q^{(m+n)/2} && \text{and} \\
  |\alpha^{-1}| &\leqslant q^{((-m-2n)+n)/2} = q^{-(m+n)/2} \text{,}
\end{align*}
hence the assertion. 





\subsection{Remark}\label{VI:1-21}

The last argument shows that if $X_0$ is a smooth separated scheme over 
$\dF_q$, if $\sF_0$ is a smooth $E_\lambda$-sheaf pointwise pure of weight 
$m$ on $X_0$, and if $\h_c^i(X,\sF) \hookrightarrow \h^i(X,\sF)$, then the 
complex conjugates of the eigenvalues of $F^\ast$ on $\h_c^i(X,\sF)$ are 
of absolute value $q^{(m+i)/2}$. 





\subsection*{Appendix. The Lang isogeny in the noncommutative case}




\subsection{}\label{VI:1-22}

Let $G_0$ be a connected algebraic group over $\dF_q$ and let $\fL$ be the morphism 
$x\mapsto F x\cdot x^{-1}$ from $G_0$ to itself. It is a morphism of 
$G_0$-homogeneous spaces if, on the source, $G_0$ acts by left translation 
$x\mapsto (y\mapsto x y)$, and on the target by 
$x\mapsto (y\mapsto F x\cdot y\cdot x^{-1})$). We have $\fL(xy) = \fL(x)$ for 
$y\in G_0(\dF_q)$, and $\fL$ induces an isomorphism $G_0/G_0(\dF_q)\iso G_0$. 

As in the commutative case, $\fL$ therefore makes $G_0$ a $G_0(\dF_q)$-torsor 
on $G_0$, the \emph{Lang torsor} $L$. If 
$\rho:G_0(\dF_q) \to \gl(n,E_\lambda)$ is a linear representation of 
$G_0(\dF_q)$, we denote by $\sF(\rho)$ the $E_\lambda$-sheaf $\rho(L)$ 
(\ref{VI:1-2}). 





\subsection{}\label{VI:1-23}

Let $\gamma\in G^F=G_0(\dF_q)$, and compute the endomorphism $F^\ast$ of 
$\sF(\rho)_\gamma$. Let $g\in G$ be such that $\fL(g)=\gamma$, and choose the frame 
$\rho(g)\in \operatorname{Isom}(E_\lambda^n,\sF(\rho)_\gamma)$. For 
$e\in E_\lambda^n$, we have 
\[
  F(\rho(g)(e)) = \rho(F g)(e) = \rho(\gamma g)(e) = \rho(g\cdot g^{-1}\gamma g)(e) \text{.}
\]
We will see that $g^{-1} \gamma g\in G_0(\dF_q)$, hence 
\[
  F(\rho(g)(e)) = \rho(g) \rho(g^{-1}\gamma g)(e)
\]
and 
\begin{equation*}\tag{1.23.1}\label{VI:eq:1-23-1}
  \rho(g)^{-1} (F_\gamma^\ast)^{-1} \rho(g) = \rho(g^{-1} \gamma g) \text{:}
\end{equation*}
$\rho(g)$ identifies the inverse of $F^\ast$ at $\gamma$ with the automorphism 
$\rho(g^{-1} \gamma g)$ of $E_\lambda^n$. It remains to understand 
$g^{-1} \gamma g$.

\begin{lemma_}\label{VI:1-24}
\begin{enumerate}[(i)]
  \item If $F g\cdot g^{-1} \in G_0(\dF_q)$, we have 
    $g^{-1} (F g\cdot g^{-1}) g = g^{-1} F g\in G_0(\dF_q)$.
  \item Let $\gamma\in G_0(\dF_q)$, and $g$ such that 
    $F g\cdot g^{-1} = \gamma$. The conjugation class of $\gamma'=g^{-1} Fg$ 
    in $G_0(\dF_q)$ only depends on that of $\gamma$. 
  \item The map induced by $\gamma\mapsto \gamma'$ on the set 
    $G_0(\dF_q)^\natural$ of conjugacy classes of $G_0(\dF_q)$ 
    is bijective. Its inverse is the corresponding map for the opposite 
    group. 
\end{enumerate}
\end{lemma_}

(i) If $\gamma= F g\cdot g^{-1}\in G_0(\dF_q)$, the identities $F g=\gamma g$ 
and 
$F\gamma=\gamma$ give $F(g^{-1} F g) = (F g)^{-1} F F g = g^{-1} \gamma^{-1} F(\gamma g) = g^{-1} \gamma^{-1} F \gamma F g = g^{-1} F g$, 
so that $g^{-1} F g \in G_0(\dF_q)$. 

(ii) For $\alpha,\gamma\in G_0(\dF_q)$, the $g$ such that 
$F g\cdot g^{-1} = \gamma$ form a right coset under $G_0(\dF_q)$, and 
if $F g\cdot g^{-1} = \gamma$, we have 
$F(\alpha g \alpha^{-1})(\alpha g \alpha^{-1})^{-1} = \alpha \gamma\alpha^{-1}$. 
As $\gamma$ ranges over a conjugacy class in $G_0(\dF_q)$, the $g$ 
such that $F g\cdot g^{-1} = \gamma$ therefore range over a double coset under 
$G_0(\dF_q)$, and $\gamma'= g^{-1} F g$ ranges over a conjugacy class. 

(iii) Finally, the inverse of $F g\cdot g^{-1}\mapsto g^{-1} F g$ is 
$g^{-1} F g\mapsto F g\cdot g^{-1}$. 





\subsection{}\label{VI:1-25}

Formula \eqref{VI:eq:1-23-1} can be read succinctly: $F^\ast$ at 
$\gamma$ is $\rho(\gamma')^{-1}$. If $\gamma$ belongs to the identity component 
of its centralizer, we can take $g$ in $Z(\gamma)$. We then have 
$\gamma=\gamma'$. If this condition is not met, $\gamma$ and $\gamma'$ 
need not be conjugate. 





\subsection{An example}\label{VI:1-26}

For $G_0=\operatorname{SL}(2)$, $q$ odd, 
$\alpha\in \dF_q^\times\setminus {\dF_q^\times}^2$, and $\bar\gamma$ is the 
conjugacy class of $-1(1+N)$, with $N^2=0$, then $\bar\gamma'$ is the 
conjugacy class of $-1\cdot (1+\alpha N)$, so that $\bar\gamma'\ne \bar\gamma$ 
if $N\ne 0$. 










\section{Dictionary}\label{VI:2}

Elementary manipulations of trigonometric sums can often 
be seen, via the trace formula, as counterparts of statements 
in cohomology. In this section, I list such statements and their 
counterparts; they carry the same number, with a $\ast$ appended to 
the cohomological statement. 





\subsection{}\label{VI:2-1}

``A sum of algebraic integers is an algebraic integer'' has as its 
cohomological analogue the following theorem, proved in 
\cite[XXI 5.2.2]{sga7} (see the use of \ref{VI:2-1}$^\ast$ in 
\ref{VI:1-13}). 





\begin{theorem*}[2.1*]\label{VI:2-1*}
Let $X_0$ be a separated scheme of finite type over $\dF_q$ and let $\sF_0$ be an 
$E_\lambda$-sheaf on $X_0$. If, for all $x\in |X_0|$, the 
eigenvalues of $F_x^\ast$ on $\sF_0$ are algebraic integers, then the 
eigenvalues of $F^\ast$ on $\h_c^i(X,\sF)$ are algebraic integers. 
\end{theorem*}





\subsection{}\label{VI:2-2}

Other integrality results are proved in 
\cite[XXI, 5.2.2 et 5.4]{sga7}: if $\dim(X_0)\leqslant n$, and $\alpha$ is 
an eigenvalue of $F^\ast$ on $\h_c^i(X,\sF)$, we have: 
\begin{enumerate}[a)]
  \item under the assumptions of \ref{VI:2-1}$^\ast$, and if $i>n$, then 
    $q^{i-n}$ divides $\alpha$;
  \item if the inverses of the eigenvalues of $F_x^\ast$ are algebraic 
    integers, then $\alpha^{-1}$ is integral away from $p$. More 
    precisely, $q^{\inf(n,i)}\alpha^{-1}$ is an algebraic integer. 
\end{enumerate}

For the sheaves considered in \ref{VI:1-9}, the eigenvalues of 
$F_x^\ast$ are roots of unity, so the assumptions of 
the theorem and of b) above are satisfied. 





\subsection{}\label{VI:2-3}

Let $f:A\to B$ be a map from one finite set to another, and 
$\varepsilon$ a function on $A$. We have 
\begin{equation*}\tag{2.3.1}\label{VI:eq:2-3-1}
  \sum_{a\in A} \varepsilon(a) = \sum_{b\in B} \sum_{f(a)=b} \varepsilon(a) \text{.}
\end{equation*}





\subsection*{2.3*}\label{VI:2-3*}

Let $f:X\to Y$ be a morphism of separated schemes of finite type over an 
algebraically closed field $k$, and let $\sF$ be an $E_\lambda$-sheaf on $X$. The 
Leray spectral sequence for $f$ in cohomology with compact support is 
\begin{align*}\tag{2.3.1*}\label{VI:eq:2-3-1*}
  E_2^{p q} = \h_c^p(Y,\eR^q f_! \sF) \Rightarrow \h_c^{p+q}(X,\sF) \text{.}
\end{align*}

Let $f_0:X_0 \to Y_0$ be a morphism of separated schemes of finite type over 
$\dF_q$ and let $\sF_0$ be an $E_\lambda$-sheaf on $X_0$. Between 
\eqref{VI:eq:2-3-1} and \eqref{VI:eq:2-3-1*} there is the following compatibility. 

\begin{enumerate}[a)]
  \item The sheaf $\eR^q f_! \sF$ is obtained by extending scalars from 
    $\dF_q$ to $\dF$ in the sheaf $\eR^q f_{0!} \sF_0$ on $Y_0$. As 
    such, it is equipped with a Frobenius correspondence 
    $F^\ast:F^\ast\eR^q f_! \sF \to \eR^q f_! \sF$. It is obtained by applying 
    the functoriality of $\eR f_!$ to the Frobenius correspondence of $\sF$: 
    namely the composite 
    \[\xymatrix{
      F^\ast \eR^q f_! \sF \ar[r]^-\sim 
        & \eR^q f_! F^\ast \sF \ar[r]^-\sim 
        & \eR^q f_! \sF \text{.}
    }\]
  \item The trace formula for $\sF_0$ reads 
    \begin{align*}\tag{1}\label{VI:eq:1}
      \sum_{x\in X_0(\dF_q)} \tr(F_x^\ast,\sF_0) = \tr(F^\ast,\h_c^\bullet(X,\sF)) \text{;}
    \end{align*}
    while that for $\eR^q f_! \sF$ reads 
    \begin{align*}\tag{2$_q$}\label{VI:eq:2_q} 
      \sum_{y\in Y_0(\dF_q)} \tr(F_y^\ast,\eR^q f_! \sF) = \tr(F^\ast, \h_c^\bullet(Y,\eR^q f_! \sF)) \text{.}
    \end{align*}
\end{enumerate}

We have, for $y\in Y(\dF)$, $(\eR^q f_! \sF)_y = \h_c^q(f^{-1}(y),\sF)$, and if 
$y\in Y_0(\dF_q)$, 
\[
  \tr(F_y^\ast,\eR^q f_! \sF) = \tr(F^\ast,\h_c^q(f^{-1}(y),\sF)) \text{,}
\]
hence by the trace formula on the fiber 
\[
  \sum_{\substack{x\in X_0(\dF_q) \\ f(x) = y}} \tr(F_x^\ast,\sF_0) = \sum (-1)^q \tr(F_y^\ast,\eR^q f_! \sF) \text{.}
\]
Taking the alternating sum of the identities \eqref{VI:eq:2_q}, we therefore find 
\begin{align*}\tag{2}\label{VI:eq:2}
  \sum_{y\in Y_0(\dF_q)} \sum_{\substack{x \in X_0(\dF_q) \\ f(x) = y}} \tr(F_x^\ast,\sF_0) = \tr(F^\ast,\h_c^\bullet(Y,\eR^\bullet f_! \sF)) \text{.}
\end{align*}

Equality of the left-hand sides of \eqref{VI:eq:1} and \eqref{VI:eq:2} 
follows from \eqref{VI:eq:2-3-1}; equality of the right-hand sides follows from 
\eqref{VI:eq:2-3-1*} (taking into account that $F^\ast$ is an endomorphism of 
the entire spectral sequence). 





\subsection{}\label{VI:2-4}

Let $(A_i)_{i\in I}$ be a finite family of finite sets, 
$A=\prod_{i\in I} A_i$, $\varepsilon_i$ a function on $A_i$, and 
$\varepsilon$ the function $a\mapsto \prod \varepsilon_i(a_i)$ on $A$. We have 
\begin{align*}\tag{2.4.1}\label{VI:eq:2-4-1}
  \sum_{a\in A} \varepsilon(a) = \prod_{i\in I} \sum_{a_i\in A_i} \varepsilon_i(a_i) \text{.}
\end{align*}





\subsection*{2.4*}\label{VI:2-4*}

Let $(X_i)_{i\in I}$ be a finite family of separated schemes of finite type 
over an algebraically closed field $k$, let $X=\prod_{i\in I} X_i$, let $\sF_i$ be an 
$E_\lambda$-sheaf on $X_i$ and $\sF$ the external tensor product 
$\bigboxtimes_{i\in I}\sF_i = \bigotimes_{i\in I} \pr_i^\ast(\sF_i)$ 
of the $\sF_i$. 
The K\"unneth formula gives 
\begin{align*}\tag{2.4.1*}\label{VI:eq:2-4-1*}
  \h_c^\bullet(X,\sF) = \bigotimes_{i\in I} \h_c^\bullet(X_i,\sF_i) \text{.}
\end{align*}

For the isomorphism \eqref{VI:eq:2-4-1*} to be canonical and independent of 
the choice of a total ordering of $I$, the tensor product on the right must be 
understood as a graded tensor product, with the Koszul sign rule 
(\hyperref[IV]{Cycle}, \ref{IV:1-3}). 





\subsection{}\label{VI:2-5}

Let $A$ be a finite set, let $B$ be a subset of $A$, and let $\varepsilon$ be a 
function on $A$. We have 
\begin{align*}\tag{2.5.1}\label{VI:2-5-1}
  \sum_{a\in A} \varepsilon(a) = \sum_{a\in A\setminus B}\varepsilon(a) + \sum_{a\in B} \varepsilon(a) \text{.}
\end{align*}





\subsection*{2.5*}\label{VI:2-5*}

Let $X$ be a separated scheme of finite type over an algebraically closed field $k$, 
let $Y$ be a closed subscheme, and let $\sF$ be an $E_\lambda$-sheaf on $X$. We have 
a long exact cohomology sequence 
\begin{align*}\tag{2.5.1*}\label{VI:eq:2-5-1*}
\xymatrix{
  \cdots \ar[r]^-\partial 
    & \h_c^i(X\setminus Y,\sF) \ar[r] 
    & \h_c^i(X,\sF) \ar[r] 
    & \h_c^i(Y,\sF) \ar[r]^-\partial 
    & \cdots 
}
\end{align*}

More generally, given a finite filtration of $X$ by closed subsets 
$X_p$ ($p\in \dZ$; assume that $X_p\supset X_{p+1}$, that $X_p=X$ 
for $p$ sufficiently small, and that $X_p=\varnothing$ for $p$ sufficiently large), there is a 
spectral sequence 
\begin{align*}\tag{2.5.2*}\label{VI:eq:2-5-2*}
  E_1^{p q} = \h_c^{p+q}(X_p\setminus X_{p+1},\sF) \Rightarrow \h_c^{p+q}(X,\sF) \text{.}
\end{align*}





\subsection{}\label{VI:2-6}

Let $A$ be a finite set, $(A_i)_{i\in I}$ a finite covering of $A$ and 
$\varepsilon$ a function on $A$. For $J\subset I$, denote by $A_J$ 
the intersection of $A_j$ ($j\in J$). We have 
\begin{align*}\tag{2.6.1}\label{VI:eq:2-6-1}
  \sum_{J\subset I} (-1)^{|J|} \sum_{a\in A_J} \varepsilon(a) = 0 \text{.}
\end{align*}





\subsection*{2.6*}\label{VI:2-6*}

Let $X$ be a separated scheme of finite type over an algebraically closed field $k$, 
$(X_i)_{i\in I}$ a closed (resp. open) finite covering of $X$ by 
subschemes, and $\sF$ an $E_\lambda$-sheaf on $X$. For $J\subset I$, 
let $X_J$ be the intersection of the $X_j$ ($j\in J$). We have the respective 
Leray spectral sequences 
\begin{align*}\tag{2.6.1*}\label{VI:eq:2-6-1*}
  E_1^{p q} &= \bigoplus_{|J|=p+1>0} \h_c^q(X_J,\sF) \otimes_\dZ \textstyle\bigwedge^{|J|} \dZ^J \Rightarrow \h_c^{p+q}(X,\sF) \\ 
  \tag{2.6.2*}\label{VI:eq:2-6-2*}
  E_1^{p q} &= \bigoplus_{|J|=1-p>0} \h_c^q(X_J,\sF) \otimes_\dZ \textstyle\bigwedge^{|J|} \dZ^J \Rightarrow \h_c^{p+q}(X,\sF)
\end{align*}
If $J=\varnothing$ were not excluded, there would also be spectral 
sequences converging to $0$. The factors $\bigwedge^{|J|} \dZ^J$ are included 
to avoid choosing a total order on $I$. 





\subsection{}\label{VI:2-7}

Let $A$ be a finite commutative group, and $\chi:A \to E_\lambda^\times$ a 
nontrivial character. We have 
\begin{align*}\tag{2.7.1}\label{VI:eq:2-7-1}
  \sum_{a\in A} \chi(a) = 0 \text{.}
\end{align*}
We will prove the cohomological analogue of \eqref{VI:eq:2-7-1} by transposing 
the following proof: if $x\in A$, $a\mapsto x a$ is a permutation of $A$, and 
\begin{align*}
  \sum_{a\in A}\chi(a) = \sum_{a\in A} \chi(x a) &= \chi(x) \sum_{a\in A} \chi(a) && \text{, whence} \\
  (\chi(x)-1)\sum_{a\in A}\chi(a) &= 0 \text{,}
\end{align*}
and by hypothesis there exists $x$ such that $\chi(x)-1\ne 0$. 





\begin{theorem*}[2.7*]\label{VI:2-7*}
Let $G_0$ be a connected commutative algebraic group over $\dF_q$ and 
$\chi:G_0(\dF_q) \to E_\lambda^\times$ a nontrivial character. We have 
\begin{align*}\tag{2.7.1*}\label{VI:eq:2-7-1*}
  \h_c^\bullet(G,\sF(\chi)) = 0 \text{.}
\end{align*}
\end{theorem*}

For $x$ a rational point of $G$, let $t_x$ be the translation $t_x(g) = x g$ 
of $G$. The formula $\fL t_x = t_{\fL(x)} \fL$ expresses that $(t_x,t_{\fL(x)})$ 
is an automorphism of the diagram $G\xrightarrow\fL G$ ($G$ equipped with the torsor of 
Lang). Let $\rho(g)$ be the automorphism of $(G,\sF(\chi))$ induced by it. 
For $g\in G_0(\dF_q)$, i.e. for $\fL(g) = e$, it is the identity on $G$, and 
multiplication by $\chi(g)^{-1}$ on $\sF(\chi)$. 

Let $\rho_H(g)$ denote the automorphism of $\h_c^\bullet(G,\sF(\chi))$ induced by 
$\rho(g)$. A homotopy argument (Lemma \ref{VI:2-8} below) shows that 
$\rho_H(g) = \rho_H(e)$, and hence is the identity. On the other hand, for 
$g\in G_0(\dF_q)$, $\rho_H(g)$ is multiplication by 
$\chi^{-1}(g)$: on $\h^\bullet(G,\sF(\chi))$, multiplication by 
$(\chi^{-1}(g)-1)$ is zero. Taking $g$ such that $\chi(g)\ne 1$, we obtain 
\ref{VI:2-7*}. 





\begin{lemma_}\label{VI:2-8}
Let $X$ and $Y$ be two schemes over an algebraically closed field $k$, with $X$ 
separated and of finite type and $Y$ connected. Let $\sF$ be a sheaf on $X$ and 
let $(\rho,\varepsilon)$ be a family of endomorphisms of $(X,\sF)$ parametrized 
by $Y$: 
\begin{align*} 
  \rho:Y\times_k X &\to Y\times_k X && \text{is a $Y$-morphism, and} \\
  \varepsilon:\rho^\ast \pr_2^\ast \sF &\to \pr_2^\ast\sF && \text{a morphism of sheaves.} 
\end{align*}
Assume that $\rho$ is proper. For $y\in Y(k)$, let $\rho_H(y)^\ast$ denote 
the endomorphism of $\h_c^\bullet(X,\sF)$ induced by $\rho_y:X\to X$ and 
$\varepsilon_y:\rho_y^\ast\sF\to \sF$. Then, $\rho_H(y)^\ast$ is 
independent of $y$. 
\end{lemma_}

Indeed, $\eR^p \pr_{1!} \pr_2^\ast\sF$ is the 
constant sheaf on $Y$ with fiber $\h_c^p(X,\sF)$, and $\rho_H(y)^\ast$ is the 
fiber at $y$ of the endomorphism 
\[\xymatrix{
  \eR^p \pr_{1!} \pr_2^\ast \sF \ar[r]^-{\rho^\ast} 
    & \eR^p \pr_{1!} \rho^\ast \pr_2^\ast \sF \ar[r]^-\varepsilon 
    & \eR^p \pr_{1!} \pr_2^\ast \sF 
}\]
of this sheaf. 





\subsection{Remark}\label{VI:2-9}

Let $G_0$ be a connected algebraic group over $\dF_q$ and 
$\rho:G_0(\dF_q) \to \operatorname{GL}(V)$ a linear representation 
such that $V^{G_0(\dF_q)}=0$. Then again $\h_c^\bullet(G,\rho(L))=0$. 










\section{Sums in one variable}\label{VI:3}





\subsection{}\label{VI:3-1}

Weil was the first to apply ``cohomological'' methods to 
the study of trigonometric sums in one variable; since he had 
proved the analogue of the Riemann hypothesis for Artin's $L$ functions 
over function fields, these methods gave him excellent 
estimates (of order $O(\text{square root of the number of terms})$) and an 
understanding of how these sums behave under ``extension of the base field.'' 

The substance of this section is a presentation, in cohomological language, of 
his results. 





\subsection{}\label{VI:3-2}

Cohomological methods here lead us to compute the groups 
$\h_c^i(X,\sF)$, where $\sF_0$ is an $E_\lambda$-sheaf on a curve $X_0$ over 
$\dF_q$. These groups vanish for $i\ne 0,1,2$, and for 
$i=0,2$, they have a simple interpretation (\ref{VI:1-8} a,b,c). Furthermore, the 
Euler-Poincaré characteristic 
$\chi_c(\sF) = \sum (-1)^i \dim \h_c^i(X,\sF)$ (which is also equal to 
$\chi(\sF) = \sum (-1)^i \dim \h^i(X,\sF)$) can be computed in terms of 
$X$, the rank of $\sF$ at the generic points of $X$, and the 
ramification of $\sF$. The essential result is the following. 

Let $\bar X$ be a smooth connected projective curve of genus $g$ over an 
algebraically closed field $k$, let $X$ be the dense open subset of $\bar X$ 
complementary to a finite set $S$ of points, and let $\sF$ be a 
smooth $E_\lambda$-sheaf on $X$. Let $\chi(X)=2-2 g-\# S$ be the 
Euler--Poincaré characteristic of $X$, and let $\operatorname{rg}(\sF)$ be the 
rank of $\sF$. For each point $s\in S$, we define an integer 
$\swan_s(\sF)$, the Swan conductor, measuring the wild ramification of 
$\sF$, and 
\begin{align*}\tag{3.2.1}\label{VI:eq:3-2-1}
  \chi_c(\sF) = \operatorname{rg}(\sF) \cdot \chi(X) - \sum_{s\in S} \swan_s(\sF) 
\end{align*}
\cite{ra65}. 





\subsection{}\label{VI:3-3}

Let $\bar X$, $j:X\hookrightarrow\bar X$ and $\sF$ as in \ref{VI:3-2}. The 
Poincaré duality theorem takes the particularly convenient form: 
$\h^i(\bar X,j_\ast \sF)$ and $\h^{2-i}(\bar X,j_\ast(\sF^\vee(1)))$ are 
dual to each other (\hyperref[V]{Duality}, \ref{V:1-3}). 





\subsection{}\label{VI:3-4}

Take $k=\dF$, and suppose that $X$, $\bar X$, and $\sF$ arise from 
$X_0$, $\bar X_0$, and $\sF_0$ over $\dF_q$. If $\sF_0$ becomes trivial on a 
finite covering of $X_0$, Weil proved that the complex conjugates of 
the eigenvalues of $F^\ast$ on $\h^i(\bar X,j_\ast\sF)$ have absolute value 
$q^{i/2}$. Such sheaves $\sF_0$ correspond to Artin $L$-functions 
over the function field of $\bar X_0$. 





\subsection*{3.5}\label{VI:3-5_} % originally one of two 3.5s

Under the same hypotheses, or more generally if $E_\lambda$ is the 
completion at a place $\lambda$ of a number field $E$ and if $\sF_0$ 
belongs to an infinite compatible system of $v$-adic representations 
($v$ a place of $E$), we can compute 
\[
  \prod_i \det(-F^\ast,\h^i(X,\sF))^{(-1)^i} 
\]
from \emph{local} information on $\sF_0$. This is explained in 
\cite{de73}. For trivial $\sF_0$ on a finite covering of $X_0$, it is 
the expression of the constant in the functional equation of an Artin 
$L$-function as a product of local constants. 





\subsection{Example}\label{VI:3-5}

Let $\psi$ be the character 
$\exp\left(\frac{2\pi i}{p} \tr_{\dF_q/\dF_p}\right)$ of $\dF_q$; we have 
$\psi(a^p-a)=1$. Let $X_0$ be a smooth projective geometrically 
irreducible curve of genus $g$ over $\dF_q$, and let $f$ be a rational function on 
$X_0$, i.e. a morphism $f:X_0 \to \dP^1$, not identically equal to 
$\infty$. We are interested in the sum 
\[
  S_f' = \sum_{\substack{x\in X_0(\dF_q) \\ f(x)\ne \infty}} \psi(f(x)) 
  \text{.}
\]
This sum is zero for $f$ of the form $g^p-g$. This suggests modifying 
the sum $S_f'$ as follows: 
\begin{enumerate}[\indent a)]
  \item for any closed point $x$ of $X_0$, set $v_x(f) = $ the order of the 
    pole of $f$ at $x$ if $f(x)=\infty$, and $v_x(f)=0$ otherwise; also set 
    $v_x^\ast(f) = \inf v_x(f+g^p-g)$ (the infimum being over $g$). 
  \item If $v_x^\ast(f)=0$, $x\in X_0(\dF_q)$, and $f+g^p-g$ is 
    regular at $x$, put $\psi(f(x)) = \psi((f+g^p-g)(x))$. Finally set 
    \begin{equation*}\tag{3.5.1}\label{VI:eq:3-5-1}
      S_f = \sum'_{x\in X_0(\dF_q)} \psi(f(x)) \text{,}
    \end{equation*}
    where $(-)'$ indicates that the sum runs over those $x$ for which 
    $v_x^\ast(f)=0$. 
\end{enumerate}

We have $S_f=S_{f+g^p-g}$. We shall verify that if $f$ is not of the 
form $g^p-g+c^{t e}$, then 
\begin{equation*}\tag{3.5.2}\label{VI:eq:3-5-2}
  |S_f|\leqslant \left(2 g-2+\sum_{v_x^\ast(f)\ne 0} [k(x):\dF_q] (1+v_x^\ast(f))\right) q^{1/2} \text{.}
\end{equation*}

We begin by passing from complex coefficients to $\ell$-adic coefficients, replacing $\psi$ by 
a character of the form $\psi_0\circ \tr_{\dF_q/\dF_p}$, with 
$\psi_0:\dF_p\hookrightarrow E_\lambda^\times$. For $j:U_0\hookrightarrow X_0$ 
the inclusion of the open where $f\ne\infty$, we then prove that 
\begin{equation*}\tag{3.5.3}\label{VI:eq:3-5-3}
  S_f = \sum (-1)^i \tr(F^\ast,\h^i(X,j_\ast \sF(\psi f))) \text{.} 
\end{equation*}

The trace formula reduces this statement to the following (for 
$x\in X_0(\dF_q)$): 
\begin{enumerate}[\indent a)]
  \item If $v_x^\ast(f)=0$, then $F_x^\ast$ on $j_\ast \sF(\psi f)$ is 
    $\psi(f(x))$; 
  \item If $v_x^\ast(f)\ne 0$, then $j_\ast \sF(\psi f)$ is ramified at $x$: 
    $(j_\ast \sF(\psi f))_{\bar x}=0$. 
\end{enumerate}

Formula a) follows from \eqref{VI:eq:1-7-6} if $v_x(f)=0$, and the general 
case reduces to this one by replacing $f$ by $f+g^p-g$: up to isomorphism, this 
does not change $\sF(\psi f)_0$ on the open where $f$ and $g$ are regular 
(\ref{VI:1-3}). 

Formula b) follows from 
\begin{equation*}\tag{3.5.4}\label{VI:eq:3-5-4}
  \swan_x \sF(f,\psi) = v_x^\ast(f) \text{.} 
\end{equation*}
After reduction to the case $v_x(f)=v_x^\ast(f)$ (hence $p\nmid v_x(f)$) 
and extension of scalars to $\dF$, this formula can be found in 
\cite[4.4]{se61}.

Finally, we check that the sheaf $\sF(\psi f)$ on $U$ is nonconstant if 
$f$ is not of the form $g^p-g+c^{t e}$: since $\psi_0$ is injective, the 
triviality of $\sF(\psi f)=\sF(\psi_0 f)$ (\ref{VI:1-8}.ii) is equivalent to 
that of the $\dF_p$-torsor on $U$ defined by $T^p-T-f=0$. If it were trivial, 
the torsor on $U_0$ defined by the same equation would be the inverse image of an 
$\dF_q$-torsor on $\spec(\dF_q)$, with equation $T^p-T-\lambda=0$, and the 
$\dF_p$-torsor equation $T^p-T-(f-\lambda)=0$ on $U_0$ would be trivial 
and $f$ would therefore be of the form $g^p-g+\lambda$. The 
$\h^i(\bar X,j_\ast \sF(\psi f))$ therefore vanish for $i\ne 1$, and formula 
\eqref{VI:eq:3-5-3} reduces to 
\begin{equation*}\tag{3.5.5}\label{VI:eq:3-5-5}
  S_f = -\tr(F^\ast,\h^1(\bar X,j_\ast \sF(\psi f))) 
\end{equation*}
and, by \eqref{VI:eq:3-2-1} and \eqref{VI:eq:3-5-4}, the dimension of $\h^1$ is 
$2 g-2+\sum_{v_x^\ast(f)>0} [k(x):\dF_q](1+v_x^\ast(f))$; moreover, $-S_f$ 
is the sum of this many eigenvalues of $F^\ast$, each having complex 
absolute value $q^{1/2}$. 





\subsection{}\label{VI:3-6}

Suppose that, for an automorphism $\sigma$ of $X_0$, we have 
\[
  f(\sigma x) = - f(x) \text{.} 
\]
The sum $S_f$ is then real. If $\sigma$ is an involution, Poincaré 
duality allows us to say a little more: set 
$\sF_0=j_\ast \sF(\psi f)$, $\sG_0 = j_\ast \sF(\psi(-f))$, we have 
\begin{enumerate}[\indent a)]
  \item $\sF_0$ and $\sG_0$ are in duality, 
  \item $\sigma^\ast \sF_0 \iso \sG_0$ and $\sigma^\ast \sG_0 \iso \sF_0$, for 
    natural isomorphisms such that the composite 
    $\sF_0=(\sigma^2)^\ast \sF_0 \iso \sigma^\ast \sG_0 \iso \sF_0$ 
    is the identity. 
  \item the pairing $\sF_0\otimes \sG_0 \to E_\lambda$ satisfies 
    $\sigma^\ast(f\cdot g) = \sigma^\ast(f) \cdot \sigma^\ast (g)$. 
\end{enumerate} 

Passing to cohomology, we find that $\h^1(X,\sF)$ and $\h^1(X,\sG)$ are in 
perfect duality with values in $E_\lambda(1)$, and that the bilinear form 
$\alpha\cdot \sigma^\ast\beta$ on $\h^1(X,\sF)$ is \emph{alternating}: 
$\sigma^\ast$ acts trivially on $E_\lambda(1)$, and 
\[
  \alpha\cdot \sigma^\ast \beta = \sigma^\ast(\alpha\cdot \sigma^\ast \beta) = \sigma^\ast \alpha\cdot \beta = -\beta\cdot \sigma^\ast \alpha \text{.} 
\]
We conclude that $\h^1(X,\sF)$ is of even dimension (we also check 
    easily that each nonzero $v_x^\ast(f)$ is odd) and that the 
eigenvalues of $F^\ast$ on $\h^1(X,\sF)$ occur in pairs $\alpha$ and 
$q/\alpha$. 





\subsection{Example}\label{VI:3-7}

This applies to the Kloosterman sums 
$\sum_{x\in \dF_q^\times} \psi\left(x+\frac a x\right)$ (take 
$X_0=\dP^1$, $f=x+\frac a x$, $\sigma x=-x$; the poles of $f$ are $0$ and 
$\infty$, and at each of them $v_x^\ast(f)=1$; the dimension of $\h^1$ is 
$-2+2+2=2$). We therefore have (for $a\ne 0$) 
\[
  \sum_{\substack{x y=a \\ x,y\in \dF_{q^n}}} \exp\left(\frac{2\pi i}{p} \tr_{\dF_{q^n}/\dF_p}(x+y)\right) = -(\alpha^n+\bar\alpha^n) \qquad \text{, }\alpha\bar \alpha = q \text{.} 
\]
This result is due to L.\ Carlitz \cite{ca69}. 

We now give an application of a Lang--Weil method. 





\begin{proposition_}\label{VI:3-8}
Let $P\in \dF_q[X_1,\dots,X_n]$ be a polynomial in $n$ variables, of degree 
$d$, which is not of the form $Q^p-Q+C^{te}$. Again setting 
$\psi(x) = \exp\left(\frac{2\pi i}{p} \tr_{\dF_q/\dF_p}(x)\right)$, we have 
\[
  \left|\sum_{x_i\in \dF_q} \psi P(x_1,\dots,x_n)\right| \leqslant (d+1) q^{n-\frac 1 2} \text{.}
\]
\end{proposition_}

For the left-hand side we have the trivial estimate $q^n$. It is therefore enough to 
prove \ref{VI:3-8} when $d-1<q^{1/2}$; assume therefore that 
$d<q+1$, and let $P_d$ be the homogeneous part of degree $d$ of $P$. Recall: 





\begin{lemma_}\label{VI:3-9}
A hypersurface of degree $d$ in $\dP^r$ ($r\geqslant 1$) cannot contain 
all rational points over $\dF_q$ unless $d\geqslant q+1$. 
\end{lemma_}

We argue by induction: if there exists a rational hyperplane not 
entirely contained in the hypersurface (this is not the case for $r=1$), we 
apply the induction hypothesis to the intersection of the hypersurface with this 
hyperplane. Otherwise, the degree $d$ is $\geqslant$ the number of rational 
hyperplanes, $\geqslant q+1$. 

Let us now distinguish two cases. 


\paragraph{Case 1: $p\nmid d$.}
Applying the lemma to $P_d$, we see that, after making a linear change 
of variables if necessary, we may assume that $P_d(1,0,\ldots,0)\ne 0$, i.e. 
that the coefficient of $X_1^d$ in $P$ is nonzero. A polynomial in $1$ 
variable 
\[
  S(X) = \sum_{i=0}^d a_i x^i \text{,} 
\]
with $a_d\ne 0$ ($p\nmid d$) is never of the form $Q^p-Q+C^{t e}$, and 
the \eqref{VI:eq:3-5-2} estimate is reduced to 
\[
  \left| \sum \psi(S(x)) \right| \leqslant (d-1) q^{1/2} \text{.} 
\]
Applying this estimate to the partial sums obtained by varying 
only $x_1$, we find 
\[
  \left| \sum \psi P(x_1,\dots,x_n) \right| = \left| \sum_{x_2,\dots,x_n} \sum_{x_1} \psi P(x_1,\dots,x_n)\right| \leqslant q^{n-1} (d-1) q^{1/2} 
\]
as promised. 


\paragraph{Case 2: $p\mid d$ ($d>0$).}
If $P_d$ is a $p$-th power, replacing $P$ by $P-(P_d-P_d^{1/p})$ 
does not change the sum under consideration and lowers the degree: this case 
is disposed of by induction on $d$. Otherwise, the differential of 
$P_d$ is not identically zero; applying \ref{VI:3-9}, one can assume, 
after a linear change of variables if necessary, that it is nonzero 
at the point $(1,0,\dots,0)$. If we write 
\[
  P_d = \sum_{i=0}^d X_1^{d-i} S_i(X_2,\dots,X_n) \text{,} 
\]
this means that the linear form $S_1$ is not identically zero. 

The form $S_0$ is constant, and by replacing $P$ with 
$P-(S_0 X_1^d -S_0^{1/p} X_1^{d/p})$, we can assume that it is zero. 
Finally, denote by $-\lambda$ the coefficient of $X_1^{d-1}$ in $P$. For fixed 
$x_2,\dots,x_n$, we have 
\[
  P(X_1,x_2,\dots,x_n) = (S_1(x_2,\dots,x_n)-\lambda) X_1^{d-1} + \text{terms of lower degree in $X_1$} 
\]
and if $S_1(x_2,\dots,x_n)\ne \lambda$, we therefore have 
\[
  \left| \sum_{x_1} \psi P(x_1,\dots,x_n)\right| \leqslant (d-2) q^{1/2} \text{.} 
\]
In total, 
\begin{align*}
  \left|\sum \psi P(x_1,\dots,x_n)\right| 
    &\leqslant \left|\sum_{S_1=\lambda} \psi P(x_1,\dots,x_n)\right| + \sum_{S_1(X_2,\dots,X_n)\ne\lambda} \left| \sum_{x_1} \psi P(x_1,\dots,x_n)\right| \\
    &\leqslant q^{n-1} + (q^{n-1}-q^{n-2})(d-2) q^{1/2} \\
    &< (d-2) q^{n-1/2} + q^{n-1} \\
    &< (d-1) q^{n-1/2} \text{.}
\end{align*}

Another result of this nature is given by R.A.\ Smith 
\cite{sm70}. 










\section{Gauss and Jacobi sums}\label{VI:4}





\subsection{}\label{VI:4-1}

Let $k$ be a finite field of characteristic $p$, let $\chi$ be a character of 
$k^\times$, and let $\psi$ be a nontrivial character of the additive group of $k$. 
We take the following as the definition of Gauss sums: 
\begin{equation*}\tag{4.1.1}\label{VI:eq:4-1-1}
  \tau(\chi,\psi) = -\sum_{x\in k^\times} \psi(x) \chi^{-1}(x) 
\end{equation*}
(note the sign). Classically, $\chi$ and $\psi$ are complex-valued. 
We take them to have values in $E_\lambda^\times$ (see the proof of 
\ref{VI:1-15}). Regard $-\tau(\chi,\psi)$ as a sum over the rational 
points of the scheme $\dG_m$ over $k$. By Section \ref{VI:1}, 
we have $-\tau(\chi,\psi) = \tr(F^\ast,\h_c^\ast(\dG_m,\sF(\psi \chi^{-1}))$. Moreover,

\begin{proposition_}\label{VI:4-2}
The cohomology of $\dG_m$ with coefficients in $\sF(\psi\chi^{-1})$ satisfies 
\begin{enumerate}[(i)]
  \item $\h_c^i=0$ for $i\ne 1$, and $\dim \h_c^1=1$. 
  \item $F^\ast$, acting on $\h_c^1$, is multiplication by 
    $\tau(\chi,\psi)$. 
  \item If $\chi$ is nontrivial, we have $\h_c^\bullet \iso \h^\bullet$. 
\end{enumerate}
\end{proposition_}
\begin{proof}[Proof]
For any $n$ prime to $p$, denote by $K_n$ the $\dmu_n$-torsor on $\dG_m$ 
defined by the exact K\"ummer sequence 
$0 \to \dmu_n \to \dG_m \xrightarrow{x^n} \dG_m \to 0$. If $k$ has $q$ 
elements, then over $k$ one has $\dmu_{q-1}=k^\times$, and the Lang torsor on 
$\dG_m/k$ is $K_{q-1}$. Hence $\sF(\chi^{-1})=\chi(K_{q-1})$. 
Assertion \ref{VI:4-2}(ii) follows from (i,iii), which in turn follow from 
the following geometric statement, where $\sF(\psi)$ denotes the sheaf 
on $\dG_m/\dF$ obtained by extending scalars from $k$ to $\dF$ in 
$\sF(\psi)$ on $\dG_m/k$ (\ref{VI:1-8}.b). 
\end{proof}





\begin{proposition_}\label{VI:4-3}
Let $\chi:\dmu_n \to E_\lambda^\times$. The cohomology of $\dG_m$ with 
coefficients in $\sF(\psi)\otimes \chi(K_n)$ satisfies 
\begin{enumerate}[\indent (i)]
  \item $\h_c^i=0$ for $i\ne 1$, and $\dim \h_c^1{} = 1$. 
  \item If $\chi$ is nontrivial, we have $\h_c^\bullet{}\iso \h^\bullet{}$. 
\end{enumerate}
\end{proposition_}
\begin{proof}[Proof]
The sheaf $\sF(\psi)$ is the restriction to $\dG_m$ of a sheaf locally 
constant on $\dG_a$, wildly ramified at infinity, with Swan conductor 
$1$. The sheaf $\chi(K_n)$ is constant if $\chi=1$; if $\chi\ne 1$, it 
is tamely ramified at $0$ and $\infty$. 

The sheaf $\sF(\psi)\otimes\chi(K_n)$ is therefore ramified at $\infty$; 
applying \ref{VI:1-18}(b,c), we find that 
$\h_c^i(\dG_m,\sF(\psi)\otimes \chi(K_n))=0$ for $i\ne 1$. If $\chi\ne 1$, it 
is ramified at $0$ and $\infty$, and (ii) follows from \ref{VI:1-19}.a. 

The Swan conductors are $0$ at $0$ and $1$ at $\infty$. By 
\eqref{VI:eq:3-2-1}, the Euler--Poincaré characteristic is therefore $-1$, 
which completes the proof. 
\end{proof}





\subsection{Remark}\label{VI:4-4}

If $\chi$ is nontrivial, \ref{VI:4-2}(iii) and Poincaré duality 
show that $\h_c^1(\dG_m,\sF(\psi\chi^{-1}))$ and 
$\h_c^1(\dG_m,\sF(\psi^{-1}\chi))$ are in duality (duality with values 
in $E_\lambda(-1)$). We therefore have 
$\tau(\chi,\psi)\cdot \tau(\chi^{-1},\psi^{-1}) = q$, i.e. 
$|\tau(\chi,\psi)|=q^{1/2}$. 





\subsection{}\label{VI:4-5}

If $k$ is an extension of degree $N$ of $\dF_q$, we can also regard 
\eqref{VI:eq:4-1-1} as a sum in $N$ variables over $\dF_q$. More 
generally, let $k$ be an étale algebra over $\dF_q$, of degree $N$ over 
$\dF_q$. It is a product of fields $k_i$ of degrees $N_i$, and we set 
\begin{equation*}\tag{4.5.1}\label{VI:eq:4-5-1}
  \varepsilon(k) = (-1)^{\sum (N_i+1)} \text{.} 
\end{equation*}
This is the sign of the permutation of $S=\hom_{\dF_q}(k,\dF)$ induced by the 
Frobenius substitution $\varphi\in \gal(\dF/\dF_q)$. 

Let $\psi:\dF_q\to E_\lambda^\times$ be a nontrivial character, and let 
$\chi:k^\times \to E_\lambda^\times$. We set 
\begin{equation*}\tag{4.5.2}\label{VI:eq:4-5-2}
  \tau_{\dF_q}(\chi,\psi) = (-1)^N \sum_{x\in k^\times} \psi \tr_{k/\dF_q}(x) \cdot \chi^{-1}(x) \text{.} 
\end{equation*}

If $\chi$ has coordinates $\chi_i:k_i^\times \to E_\lambda^\times$, we 
have the elementary identity 
\begin{equation*}\tag{4.5.3}\label{VI:eq:4-5-3}
  \tau_{\dF_q}(\chi,\psi) = \varepsilon(k) \prod \tau(\chi_i,\psi\circ \tr_{k_i/\dF_q}) \text{.}
\end{equation*}

After some preliminaries, we give in \ref{VI:4-10} a 
cohomological interpretation of the sums \eqref{VI:eq:4-5-2}. In 
\ref{VI:4-12}, we will interpret the Hasse--Davenport identity as a 
twisted form (cf. \ref{VI:1-12}) of the following special case of 
\eqref{VI:eq:4-5-3}: for $k=\dF_q^N$, and $\chi$ a character of 
$\dF_q^\times$, we have 
$\tau_{\dF_q}(\chi\circ N_{k/\dF_q},\psi)=\tau(\chi,\psi)^N$. 





\subsection{}\label{VI:4-6}

Recall that for $M$ a finitely generated projective module over a ring $A$, the 
functor $\spec(B)\mapsto M\otimes_A B$ from affine schemes over $\spec(A)$ 
to $\mathsf{Ens}$ is represented by the affine scheme $\dV(M^\vee)$ 
(EGA notation), the spectrum of $\operatorname{Sym}_A^\bullet(M^\vee)$. If 
$M=A^n$, this is the standard affine space of dimension $n$. 

Suppose that $M$ is a unital $A$-algebra, and let 
$V=\dV(M^\vee)$. By definition, for any $B$ equipped with an $A$-algebra structure, the set 
$V(B)$ of points of $V$ with coordinates in $B$ is $M\otimes_A B$. Since the 
functor $B\mapsto M\otimes_A B$ takes values in unital rings, 
$V$ is a unital ring scheme over $\spec(A)$. The 
norm and trace morphisms $M\otimes_A B\to B$ are functorial in $B$; they 
therefore correspond to morphisms of schemes $N$ and $T$ from $V$ to 
$\dG_a$. We shall consider the following schemes obtained from $V$: 
\begin{enumerate}[a)]
  \item $V^\ast$ is the open subscheme of invertible elements of $V$ (a group 
    scheme under $\cdot$);
  \item $W$ is the hyperplane of equation $T=0$ and $W^\ast=W\cap V^\ast$; 
  \item $P$ is the hyperplane at infinity $V\setminus \{0\}/\dG_m$ of the affine 
    space $V$ over $\spec(A)$. If $Q$ is the hyperplane at infinity of $W$, 
    $W^\ast/\dG_m$ and $V^\ast/\dG_m$ are open subsets of the projective spaces 
    $Q$ and $P$ over $\spec(A)$. 
    \begin{equation*}\tag{4.6.1}\label{VI:eq:4-6-1}
    \xymatrix{
      W^\ast \ar@{^{(}->}[r] \ar[d]^-\pi 
        & V^\ast \ar[d]^-\pi \\
      W^\ast/\dG_m \ar@{^{(}->}[r] 
        & V^\ast/\dG_m 
    }\qquad\subset\qquad
    \xymatrix{
      W\setminus \{0\} \ar@{^{(}->}[r] \ar[d]^-\pi 
        & V\setminus \{0\} \ar[d]^-\pi \\
      Q \ar@{^{(}->}[r] 
        & P \text{.}
    }
    \end{equation*}
\end{enumerate}

If $M=A^I$, $I$ a finite set, then $V\simeq \dG_a^I$, $V^\ast=\dG_m^I$, 
$V^\ast/\dG_m$ is the torus $\dG_m^I/(\dG_m\text{ diagonal})$, $N$ and $T$ 
are given by $\prod x_i$ and $\sum x_i$, and $Q$ is therefore the projective 
hyperplane with equation $\sum x_i=0$ in $P$. 

The case of interest is that in which $M$ is finite étale over $A$. The 
previous description then applies locally to $\spec(A)$ (for the 
étale topology) and we can write $V^\ast=\dG_m^I$, for $I$ a locally 
constant sheaf of finite sets on $\spec(A)$. If, for example, $A$ is a 
field with algebraic closure $\bar A$, and $M$ is a separable $A$-algebra, 
we put $I=\hom_A(M,\bar A)$, and over $\bar A$ we have 
$V\sim \dG_a^I$. Via this isomorphism, $N$ and $T$ are given by $\prod x_i$ and 
$\sum x_i$. 





\subsection{The Kummer torsor}\label{VI:4-7}

Let $S$ be a scheme, let $n$ be an integer invertible on $S$, and let $G$ be a 
commutative group scheme with connected fibers (written multiplicatively) over $S$. 
The sequence 
\[\xymatrix{
  0 \ar[r] 
    & G_n \ar[r] 
    & G \ar[r]^-{x^n} 
    & G \ar[r] 
    & 0 
}\]
is exact. It defines a $G_n$-torsor $K_n(G)$ (or simply $K_n$) on 
$G$. For $\chi:G_n \to E_\lambda^\times$ a homomorphism from the étale 
$S$-sheaf $G_n$ to the constant sheaf $E_\lambda^\times$, we denote by 
$\sK_n(\chi)$ the $E_\lambda$-sheaf $\chi^{-1}(K_n)$. 

The $K_n$ torsors form an inverse system of torsors under the 
projective system of groups $G_n$ (transition morphisms 
$x\mapsto x^d:G_{n d}\to G_n$). This expresses the commutativity of the diagrams 
\[\xymatrix{
  0 \ar[r] 
    & G_{n d} \ar[r] \ar[d]^-{x^d} 
    & G \ar[r]^-{x^{n d}} \ar[d]^-{x^d} 
    & G \ar[r] \ar@{=}[d] 
    & 0 \\
  0 \ar[r] 
    & G_n \ar[r] 
    & G \ar[r]^-{x^n} 
    & G \ar[r] 
    & 0 \text{.} 
}\]
We therefore have $\sK_n(\chi) = \sK_{n d}(\chi\circ x^d)$. 





\subsection{}\label{VI:4-8}

Apply the construction of \ref{VI:4-6} with $A=\dF_q$ and with $M=k$ an 
étale algebra over $\dF_q$. In accordance with our general conventions, we denote 
by a subscript $0$ the $\dF_q$-schemes denoted $V,V^\ast,\ldots$ in 
\ref{VI:4-6}. The same letters without a subscript denote the schemes over 
$\dF$ obtained by extension of scalars. 

We put $I=\hom(k,\dF)$; then $V\sim \dG_a^I$ and $T$ is given by 
$(x_i)\mapsto \sum x_i$. 
\begin{equation*}\tag{4.8.1}\label{VI:eq:4-8-1}
  \sF(\psi\circ \tr_{k/\dF_q}) = T^\ast \sF(\psi) = \bigotimes_i \pr_i^\ast \sF(\psi) \text{.} 
\end{equation*}
Since $V^\ast\sim \dG_m^I$, a character $\chi$ of $V_n^\ast\sim \dmu_n^I$, 
with values in $E_\lambda^\times$, amounts to a family 
$(\chi_i)_{i\in I}$ of characters of $\dmu_n$ indexed by $I$. It is 
defined over $\dF_q$ if, for all $\sigma\in \gal(\dF/\dF_q)$, we have 
$\chi_{\sigma_i} = \chi_i\circ \sigma^{-1}$. It then defines a 
$E_\lambda$-sheaf $\sK_n((\chi_i)_{i\in I})$ on $V_0^\ast$. If the product 
of the $\chi_i$ is trivial, $\chi$ factors through a character of 
$(V_0^\ast/\dG_m)_n$ and $\sK_n((\chi_i)_{i\in I})$ is the inverse image 
of a sheaf, denoted in the same way, on $V_0^\ast/\dG_m$. This construction 
is related to the Lang torsor as follows. 





\begin{lemma_}\label{VI:4-9}
\begin{enumerate}[(i)]
  \item For $n$ sufficiently divisible, we have a commutative diagram 
    \[\xymatrix{
      0 \ar[r] 
        & (V_0^\ast)_n \ar[r] \ar[d]^-\tau 
        & V_0^\ast \ar[r]^-{x^n} \ar[d]^-\tau 
        & V_0^\ast \ar[r] \ar@{=}[d] 
        & 0 \\
      0 \ar[r] 
        & k^\times \ar[r] 
        & V_0^\ast \ar[r]^-\fL 
        & V_0^\ast \ar[r] 
        & 0 \text{;} 
    }\]
    if $\chi$ is a character $\chi:k^\times \to E_\lambda^\times$, we have 
    $\sF(\chi) = \sK_n(\chi\circ\tau)$. 
  \item For sufficiently divisible $n$, $\tau$ identifies $k^\times$ with the coinvariants 
    of $\gal(\dF/\dF_q)$ acting on $V_n^\ast\sim \dmu_n^I$. 
  \item For $k$ a field, $N=[k:\dF_q]$, $\omega\in I$ an embedding of $k$ 
    in $\dF$, and $n=q^N-1$, the index component $\omega$ of 
    $\chi\circ\tau$ is $\chi\circ \omega^{-1}$. 
  \item If $\chi$ is nontrivial on each factor of $k$, the 
    $(\chi\circ \tau)_i$ are all nontrivial. 
\end{enumerate}
\end{lemma_}

It is enough to prove the lemma when $k$ is a field of degree $N$ over 
$\dF_q$. In this case, $n$ is ``sufficiently divisible'' if $q^N-1\mid n$. 
Choose an embedding $\omega$ of $k$ into $\dF$; then $I$ is identified 
with $\dZ/N$: the element $i\in \dZ/N$ corresponds to $\omega_i=\omega^{q^i}$. Via 
the isomorphism $V_0^\ast(\dF) = {\dF^\times}^I$, we have 
\begin{enumerate}[\indent a)]
  \item $x\in k^\times$ corresponds to $(\omega_i(x))\in {\dF^\times}^I$; 
  \item $F((x_i)_{i\in \dZ/N})=(x_{i-1}^q)_{i\in \dZ/N}$. 
\end{enumerate}

A character $\chi=(\chi_i)$ of $V_n^\ast\sim \dmu_n^I$ is defined over 
$\dF_q$ if $\chi_{-i}(x) =\chi_0(x^{q^i})$ ($i\in \dZ$). If $q^N-1\mid n$, there are 
$q^N-1$ such characters: $\chi_0$ factors through $\dmu_{q^N-1}$ and 
determines the $\chi_i$. To prove (ii), it is therefore sufficient to verify (i) 
and (iii) (or its corollary (iv)) which ensures that $\chi\mapsto \chi\circ \tau$ 
is injective. 

Let us prove (i) and (iii) for $n=q^N-1$. Write the group of 
endomorphisms of $V^\ast\sim \dG_m^I$ additively; in particular, denote by $n$ 
the $x\mapsto x^n$ operator. If $\alpha$ is the cyclic permutation operator 
$(x_i)\mapsto (x_{i-1})$, we have 
$q^N-1=(q\alpha)^N-1 = (q\alpha-1)((q\alpha)^{N-1}+\cdots + 1)$. This 
determines $\tau$: we have $\tau((x_i)) = \prod_{0\leqslant j<N} x_{i-j}^{q^j}$, 
and the induced map from $(V_0^\ast)_n$ to $k^\times$ is [as in a) above] 
$(x_i)\mapsto \prod x_{-i}^{q^i}$; hence (iii) follows. 





\begin{proposition_}\label{VI:4-10}
The cohomology of $V^\ast$ with coefficients in 
$\sF(\chi^{-1}\cdot \psi\tr_{k/\dF_q})$ satisfies 
\begin{enumerate}[\indent (i)]
  \item $\h_c^i = 0$ for $i\ne N$, and $\dim \h_c^N=1$. 
  \item On $\h_c^N$, $F^\ast$ is multiplication by 
    $\tau_{\dF_q}(\chi,\psi)$. 
  \item If $\chi$ is nontrivial on each factor of $k$, we have 
    $\h_c^\bullet \iso \h^\bullet$. 
\end{enumerate}
\end{proposition_}

Applying \ref{VI:4-9}(i), over $\dF$, if $\chi\circ \tau=(\chi_i)_{i\in I}$, 
we have $\sF(\chi) = \sK_n((\chi_i)_{i\in I})$. Points (i) and (iii) 
therefore follow from the following more geometric statement (for (iii), 
apply \ref{VI:4-9}(iv)), and (ii) follows from the trace formula: 





\begin{proposition_}\label{VI:4-11}
Let $(\chi_i)_{i\in I}$ be a family of characters of $\dmu_n(\dF)$. The 
cohomology of $V^\ast\simeq \dG_m^I$ with coefficients in 
$\sK_n((\chi_i)_{i\in I})\otimes \sF(\psi\tr_{k/\dF_q})$ satisfies 
\begin{enumerate}[\indent (i)]
  \item $\h_c^i{}=0$ for $i\ne N$, and $\dim \h_c^N{} = 1$. 
  \item If the $\chi_i$ are all nontrivial, we have 
    $\h_c^\bullet{}\iso \h^\bullet$. 
\end{enumerate}
\end{proposition_}

We have $V^\ast\sim \dG_m^I$, 
$\sK((\chi_i)) = \bigotimes_i \pr_i^\ast \chi_i(K_n(\dG_m))$ and 
$\sF(\psi\circ \tr_{k/\dF_q})=T^\ast \sF(\psi) = \bigotimes_i \pr_i^\ast \sF(\psi)$. 
The K\"unneth formula then applies, and \ref{VI:4-11} follows from 
\ref{VI:4-3}. 





\subsection{}\label{VI:4-12}

From \ref{VI:2-4*}* and (\hyperref[IV]{Cycle}, \ref{IV:1-3} Example 2), it also follows 
that the group of permutations $\sigma$ of $I$ such that 
$\chi_i = \chi_{\sigma i}$ ($i\in I$) acts on $\h_c^N$ by multiplication by 
the sign $\varepsilon(\sigma)$. We deduce a second proof 
of the identity of Hasse-Davenport. 

If $\chi$ is a character of $\dF_q^\times$, the fact that, over $\dF$, $N$ 
is given by $(x_i)\mapsto \prod x_i$ gives 
$\sF(\chi\circ N)=N^\ast\sF(\chi) = \bigotimes_i \pr_i^\ast\sF(\chi)$ 
and the K\"unneth formula gives 
\[
  \h_c^\bullet(V^\ast,\sF(\chi^{-1}\circ N,\psi\circ \tr)) \sim \h_c^\bullet(\dG_m,\sF(\chi^{-1}\psi))^{\otimes I} \text{,} 
\]
where the tensor product on the right-hand side is taken in the sense of 
\ref{VI:2-4*}*. Since $\h_c^1(\dG_m,\sF(\psi\chi^{-1}))$ is of dimension 
$1$, its tensor power $\otimes I$, in the ordinary sense, depends only on 
the cardinality $n$ of $I$. Applying (\hyperref[IV]{Cycle}, \ref{IV:1-3} Example 
2), we obtain a canonical isomorphism 
\begin{equation*}\tag{4.12.1}\label{VI:eq:4-12-1}
  \h_c^N\left(V^\ast,\sF(\chi^{-1}\circ N)\otimes T^\ast \sF(\psi)\right) \sim \h_c^1(\dG_m,\sF(\psi\chi^{-1})) \otimes_\dZ \textstyle \bigwedge^N \dZ^I \text{.}
\end{equation*}
To compute the action of $F^\ast$, it is most convenient to adopt the Galois 
point of view and observe that, since the isomorphism \eqref{VI:eq:4-12-1} is 
canonical, it is compatible with the transport-of-structure action of 
$\gal(\dF/\dF_q)$. If $\varepsilon(k)$ is the sign of the permutation 
$\varphi$ of $I$, we find that 
\begin{align*} 
  \tau_{\dF_q}(\chi\circ N_{k/\dF_q},\psi) 
    &= \tr\left(\varphi^{-1},\h_c^N\left(V^\ast,N^\ast\sF(\chi^{-1})\otimes T^\ast \sF(\psi)\right)\right) \\ 
    &= \varepsilon(k) \tr\left(\varphi^{-1},\h_c^1(\dG_m,\sF(\chi^{-1}\psi))\right)^N \\
    &= \varepsilon(k) \tau(\chi,\psi)^N \text{,} 
\end{align*}
that is, 
\begin{equation*}\tag{4.12.2}\label{VI:eq:4-12-2}
  \tau_{\dF_q}\left(\chi\circ N_{k/\dF_q},\psi\right) = \varepsilon(k) \tau(\chi,\psi)^N \text{.} 
\end{equation*}

If $k$ is a field, $\varphi$ is a cyclic permutation of $I$, 
$\varepsilon(k) = (-1)^{N+1}$, and we obtain the Hasse--Davenport identity: 
\[
  \tau\left(\chi\circ N_{k/\dF_q},\psi\circ \tr_{k/\dF_q}\right) = (-1)^{N+1} \tau_{\dF_q}\left(\chi\circ N_{k/\dF_q},\psi\right) = \tau(\chi,\psi)^N \text{.} 
\]





\begin{lemma_}\label{VI:4-13}
For $\chi\circ \tau=(\chi_i)_{i\in I}$ as in \ref{VI:4-9}(i), the
following conditions are equivalent 
\begin{enumerate}[\indent (i)]
  \item $\chi| \dF_q^\times$ is trivial. 
  \item The product of $\chi_i$ is trivial. 
  \item $\sF(\chi) = \sK_n((\chi_i))$ is the inverse image of a unique 
    sheaf on $V^\ast/\dG_m$. 
  \item The inverse image of $\sF(\chi)$ on $\dG_m$ (embedded in 
    $V^\ast$ by the structural morphism $\dF_q \to k$) is trivial. 
\end{enumerate}
\end{lemma_}

(i)$\Rightarrow$(ii). We regard $\chi$ as a character of 
$V^\ast/\dG_m(\dF_q)=k^\times/\dF_q^\times$, and we take $\sF(\chi)$ on 
$V^\ast/\dG_m$. 

(ii)$\Rightarrow$(iii). Likewise, we regard $(\chi_i)$ as a character 
of $(V^\ast/\dG_m)_n$. The uniqueness in (iii) follows from the fact that $V^\ast$ is 
a fibration with connected fibers over $V^\ast/\dG_m$, and (iii)$\Rightarrow$(iv)  
is trivial. 

(iv)$\Rightarrow$(i), (ii). This inverse image is 
$\sF(\chi|\dF_q^\times)$ and $\sK_n(\prod \chi_i)$. 





\subsection{}\label{VI:4-14}

Let $k$ be an étale algebra over $\dF_q$, of 
dimension $N+1$, and let $\chi$ be a nontrivial character of $k^\times$, trivial on 
$\dF_q^\times$. The Jacobi sum $J(\chi)$ is defined by 
\begin{equation*}\tag{4.14.1}\label{VI:eq:4-14-1}
  J(\chi) = (-1)^{N-1} \sum_{\substack{x\in k^\times/\dF_q^\times \\ \tr(x)=0}} \chi^{-1}(x) \text{.} 
\end{equation*}

We have the following identity between Gauss sums and Jacobi sums. 





\begin{proposition_}\label{VI:4-15}
For $\chi$ as above and $\psi$ a nontrivial additive character of 
$\dF_q$, we have 
\[
  q J(\chi) = \tau_{\dF_q}(\chi,\psi) \text{.} 
\]
\end{proposition_}

In the special case $k=\dF_q^n$, formula \eqref{VI:eq:4-5-3} rewrites this as 
\begin{equation*}\tag{4.15.1}\label{VI:eq:4-15-1}
  q J(\chi) = \prod_i \tau(\chi_i,\psi) \text{,} 
\end{equation*}
or equivalently as 
\begin{align*}
  \chi_0(-1) J(\chi) &= (-1)^{N-1} \sum_{\substack{x_1,\dots,x_N\in \dF_q^\times \\ \sum x_i = 1}} \prod_{i=1}^N \chi_i^{-1}(x_i) = \tau(\chi_0^{-1},\psi)^{-1} \prod_{i=1}^N \tau(\chi_i,\psi) \\
  \chi_0^{-1} &= \prod_{i=1}^N \chi_i 
\end{align*}





\begin{proof}[Proof of \ref{VI:4-15}]
\[
  \tau_{\dF_q}(\chi,\psi) = (-1)^{N+1} \sum_{x\in k^\times} \chi(x)^{-1} \psi \tr(x) = (-1)^{N+1} \sum_{x\in k^\times/\dF_q^\times} \chi(x)^{-1} \sum_{\lambda\in \dF_q^\times} \psi(\lambda \tr x) \text{.}
\]
The sum $\sum_{\lambda\in \dF_q^\times} \psi(\lambda \tr x)$ equals $q-1$ if 
$\tr(x)=0$, and $-1$ if $\tr(x)\ne 0$. Therefore, 
\[
  \tau_{\dF_q}(\chi,\psi) = (-1)^{N+1} \left(q\sum_{\substack{x\in k^\times/\dF_q^\times \\ \tr(x)=0}} \chi(x)^{-1} - \sum_{x\in k^\times/\dF_q^\times} \chi(x)^{-1}\right) \text{.}
\]
The second term on the right-hand side is zero, since it is the sum of the values 
of a character, and \ref{VI:4-15} follows. 
\end{proof}

Proposition \ref{VI:4-15} therefore has the following cohomological counterpart: 





\begin{proposition_}\label{VI:4-16}
The cohomology of $W^\ast/\dG_m$ (\ref{VI:4-6}, \ref{VI:4-8}) with coefficients 
in $\sF(\chi^{-1})$ satisfies 
\begin{enumerate}[\indent (i)]
  \item $\h_c^i=0$ for $i\ne N-1$, and $\dim \h_c^{N-1}{}=1$. 
  \item On $\h_c^{N-1}$, $F^\ast$ is multiplication by $J(\chi)$. 
  \item $\h_c^{N-1}$ is canonically isomorphic to 
    $\h_c^{N+1}\left(V^\ast,\sF(\chi^{-1},\psi\circ \tr_{k/\dF_q})\right)(1)$. 
\end{enumerate}
\end{proposition_}

Assertion (ii) follows from (i) and the trace formula; (i) and (iii) 
follow from \ref{VI:4-11} and the following more geometric statement. 





\begin{proposition_}\label{VI:4-17}
Let $(\chi_i)_{i\in I}$ be a family of characters of 
$\dmu_n(\dF)$ that are not all trivial and whose product is $1$. Then 
$\h_c^{i-1}\left(W^\ast/\dG_m,\sK_n((\chi_i)_{i\in I})\right)$ is 
canonically isomorphic to 
\newline
$\h_c^{i+1}\left(V^\ast,\sK_n((\chi_i)_{i\in I})\otimes \sF(\psi\tr_{k/\dF_q})\right)(1)$. 
\end{proposition_}

The first line of \ref{VI:4-15}, namely the projection formula, gives ($\pi$ as in 
\eqref{VI:eq:4-6-1}). 
\begin{equation*}\tag{4.17.1}\label{VI:eq:4-17-1}
  \eR^i \pi_!\left(\sK_n((\chi_i)_{i\in I})\otimes \sF(\psi \tr_{k/\dF_q})\right) = \sK_n((\chi_i)_{i\in I}) \otimes \eR^i \pi_! \sF(\psi\tr_{k/\dF_q}) 
\end{equation*}
(on $V^\ast/\dG_m$). Let us compute the sheaves 
$\eR^i\pi_! \sF(\psi\tr_{k/\dF_q}) = \eR^i \pi_! T^\ast \sF(\psi)$. Let 
$v_0:\widetilde V_0 \to V_0$ be the blow-up of $V_0$ at $\{0\}$. In the 
diagram 
\begin{equation*}\tag{4.17.2}\label{VI:eq:4-17-2}
\xymatrix{
  V_0\setminus \{0\} \ar@{^{(}->}[r] \ar[dr]_-\pi 
    & \widetilde V_0 \ar[r]^-{v_0} \ar[d]^-{\bar\pi} 
    & V_0 \\
  & P_0 
}
\end{equation*}
$\pi$ is a fibration whose fibers are punctured affine lines, $\widetilde V_0$ 
is the corresponding line bundle, and $Z_0=v_0^{-1}(0)$ is its section $0$. 
The $T^\ast\sF(\psi)$ sheaf on $V_0\setminus \{0\}$ extends to 
$(T v_0)^\ast \sF(\psi)$ on $\widetilde V_0$, and the restriction of this sheaf 
to $Z_0$ is the constant sheaf $E_\lambda$, because $T_0 v_0$ vanishes on 
$Z_0$. The long exact sequence for cohomology with compact supports is therefore 
\begin{equation*}\tag{4.17.3}\label{VI:eq:4-17-3}
\xymatrix{
  \ar[r] 
    & \eR^i \pi_! T^\ast \sF(\psi) \ar[r] 
    & \eR^i \bar\pi_! (T v_0)^\ast \sF(\psi) \ar[r] 
    & (E_\lambda\text{ for $i=0$, $0$ otherwise}) \ar[r] 
    & 
}
\end{equation*}

The inverse image $\bar\pi^{-1}(Q_0)$ is the blow-up $\widetilde W_0$ of 
$W_0$ at $0$; $T_0 v_0$ vanishes on $\widetilde W_0$; hence 
$(T v_0)^\ast \sF(\psi)=E_\lambda$ on $\bar\pi^{-1}(Q_0)$ and, since $\bar\pi$ is an 
affine-line fibration, 
\begin{equation*}\tag{4.17.4}\label{VI:eq:4-17-4}
  \eR^i \bar\pi_! (T v_0)^\ast \sF(\psi)|Q = 
    \begin{cases}
      0 & \text{for $i\ne 2$} \\
      E_\lambda(-1) & \text{for $i=2$.} 
    \end{cases} 
\end{equation*}

If $x\in P$, $x\notin Q$, the line $D=\pi^{-1}(x)$ is mapped 
isomorphically onto $\dG_a$ by $T v_0$; therefore (by 
\ref{VI:2-7}*) 
\begin{align} \notag 
  \h_c^\bullet(D,(T v_0)^\ast \sF(\psi)) &= \h_c^\bullet(\dG_a,\sF(\psi)) = 0 && \text{and} \\
  \tag{4.17.5}\label{VI:eq:4-17-5}
  \eR^i \bar\pi_! (T v_0)^\ast \sF(\psi) 
    &= \begin{cases} 
         0 & \text{for $i\ne 2$} \\
         E_\lambda(-1)_Q & \text{for $i=2$} 
       \end{cases} 
\end{align}

Combining \eqref{VI:eq:4-17-3} and \eqref{VI:eq:4-17-5}, we finally find 
\begin{equation*}\tag{4.17.6}\label{VI:eq:4-17-6}
  \eR^i \pi_! T^\ast\sF(\psi) = 
    \begin{cases}
      0 & \text{for $i\ne 1,2$} \\
      E_\lambda & \text{for $i=1$} \\
      E_\lambda(-1)_Q & \text{for $i=2$.} 
    \end{cases}
\end{equation*}





\subsection{}\label{VI:4-18}

Let us compute the compactly supported cohomology of the sheaf 
$\sK_n((\chi_i))\otimes T^\ast \sF(\psi)$ on $V^\ast$ using the 
Leray spectral sequence of $\pi:V^\ast\to V^\ast/\dG_m$. Applying 
\eqref{VI:eq:4-17-1} and \eqref{VI:eq:4-17-6}, we obtain as initial terms 
$E_2^{p,1} = \h_c^p(V^\ast/\dG_m,\sK_n(\chi_i)) = 0$ (because $(\chi_i)\ne 1$) 
and $E_2^{p,2} = \h_c^p(W^\ast/\dG_m,\sK(\chi_i))(-1)$. 

The spectral sequence reduces to an isomorphism, and \ref{VI:4-17} follows. 





\subsection{}\label{VI:4-19}

The sheaves $\sK_n((\chi_i)_{i\in I})$ make sense over any field, 
or even over any base scheme. This allows us to generalize 
\ref{VI:4-16}(i). With the notation of \ref{VI:4-6}, suppose that $M$ is finite 
étale of constant rank $N$ over $A$, and let $n$ be invertible in $A$. 
Let $a$ be the projection of $W^\ast/\dG_m$ to $\spec(A)$. Let also $\chi$ be 
a character (morphism of sheaves) 
$\chi:(V^\ast/\dG_m)_n \to E_\lambda^\times$, and $\sK_n(\chi)$ the sheaf 
corresponding. Locally for the étale topology, we can regard $\chi$ 
as a family of characters $(\chi_i)_{i\in I}$ of $\dmu_n$, whose product is 
trivial.

\begin{proposition_}\label{VI:4-20}
\begin{enumerate}[(i)]
  \item If $\chi$ is everywhere nontrivial, we have 
    $\eR^i a_! (\sK_n(\chi)) = 0$ for $i\ne N-2$, and 
    $\eR^{N-2} a_!(\sK_n(\chi))$ is smooth of rank $1$. 
  \item An automorphism $\sigma$ of $M$ which preserves $\chi$ acts on this 
    $\eR^{N-2} a_!$ by multiplication by the sign $\varepsilon(\sigma)$ 
    of $\sigma$, seen as a permutation of $I$. 
  \item If the $\chi_i$ are all nontrivial, we have $\eR a_! \iso \eR a_\ast$. 
\end{enumerate}
\end{proposition_}
\begin{proof}[Proof]
The scheme $W^\ast/\dG_m$ is the complement of a relative normal-crossings 
divisor in the scheme $Q$, which is proper and smooth over $\spec(A)$, and 
$\sK_n(\chi)$ is locally constant on $W^\ast/\dG_m$ and is tamely ramified 
along the boundary. Consequently, the $\eR^i a_!$ and 
$\eR^i a_\ast$ are smooth sheaves, compatible with every change of 
base. A standard argument then reduces us to the case in which $A$ is a 
finite field, and (i) follows from \ref{VI:4-17} and \ref{VI:4-11}(i). 
To prove (ii), use the compatibility of the action of $\sigma$ with 
the isomorphism \ref{VI:4-17}, together with \ref{VI:4-12}. For (iii), note that if the 
$\chi_i$ are all nontrivial, then $\sK_n(\chi)$ on 
$W^\ast/\dG_m\subset Q$ is ramified along every divisor at infinity, 
and apply \ref{VI:1-19-1}. 
\end{proof}










\section{Hecke characters}\label{VI:5}





\subsection{}\label{VI:5-1}

Let $F$ be a number field (of finite degree over $\dQ$) and let $k$ be a field of 
characteristic $0$. Here are several equivalent descriptions of an 
\emph{algebraic} homomorphism $\alpha:F^\times \to k^\times$. 
\begin{enumerate}[(i)]
  \item Let $\{e_i\}$ be a basis of $F$ over $\dQ$. The homomorphism $\alpha$ is 
    algebraic if it is given by a formula 
    $\alpha(\sum x^i e_i) = A(x^i)$, $A\in k(X^i)$. 
\end{enumerate}

This means that $\alpha$ agrees on $F^\times$ with a rational map defined over $k$ 
from $\res_{F/\dQ}(\dG_m)$ to $\dG_m$. By the Zariski density of 
$F^\times$, and because a rational homomorphism is defined everywhere, such 
a map is a homomorphism of group schemes: 
\begin{enumerate}[(i)]
\setcounter{enumi}{1}
  \item $\alpha$ is induced by a homomorphism of group schemes over $k$ 
    \[
      \res_{F/\dQ}(\dG_m)\otimes_\dQ k \to \dG_m \text{.} 
    \] 
\end{enumerate}

If $k$ is a number field, the adjunction property of restriction of 
scalars $R$ shows that this is equivalent to 
\begin{enumerate}[(i)]
\setcounter{enumi}{2}
  \item ($k$ a number field) $\alpha$ is induced by a homomorphism of 
    group schemes over $\dQ$: $\res_{F/\dQ}(\dG_m) \to \res_{k/\dQ}(\dG_m)$. 
\end{enumerate}

Let $\bar k$ be an algebraic closure of $k$, and let $I=\hom(F,\bar k)$. Over 
$\bar k$, the character group of $\res_{F/\dQ}(\dG_m)$ is $\dZ^I$, with basis 
the embeddings of $F$ into $\bar k$. The characters defined over $k$ 
are those invariant under $\gal(\bar k/k)$; they can be described either as 
those sending $F^\times$ into $k^\times\subset \bar k$, or in terms of 
orbits of $\gal(\bar k/k)$ on $I$, which in turn correspond to the factors of 
$F\otimes k$. 
\begin{enumerate}[(i)]
\setcounter{enumi}{3}
  \item $\alpha$ is of the form 
    $\alpha=\prod_{\omega\in I} \omega^{n_\omega}$. The allowed families of exponents 
    $(n_\omega)$ are those such that $n_\omega=n_{\omega'}$ for 
    $\omega$ and $\omega'$ in the same orbit of $\gal(\bar k/k)$. Equivalently, these 
    are those such that $\prod_\omega \omega(x)^{n_\omega}\in k$ for every 
    $x\in F$. 
  \item Write $F\otimes k=\prod_{j\in J} F_j$, where the $F_j$ are fields. 
    Then $\alpha$ can be written $\alpha=\prod N_{F_j/k}^{m_j}$. 
\end{enumerate}

In the special case where $k$ contains a normal closure of $F$, these 
expressions become: $\alpha=\prod \omega^{n_\omega}$, where $\omega$ 
runs over the $[F:\dQ]$ embeddings of $F$ into $k$. 





\subsection{}\label{VI:5-2}

Let $k$ be a number field, and let $\alpha$ be an algebraic homomorphism 
from $F^\times$ to $k^\times$. There is then a unique 
homomorphism, also denoted $\alpha$, of the group $I(F)$ 
of fractional ideals of $F$ to that of $k$, such that $\alpha((x))=(\alpha(x))$. 
Uniqueness follows from the fact that every ideal $a$ has a power that is 
principal, and that $I(k)$ is torsion-free. To prove 
existence, we use for example \ref{VI:5-1}(v): we have 
$\alpha(a) = \prod N_{F_j/k}^{m_j}((a))$. 





\subsection{}\label{VI:5-3}

Let us recall the definition of algebraic Hecke characters, also called 
Weil Hecke characters (or gr\"ossencharaktere) of type $A_0$. Let $F$ be 
a number field, let $\fm$ be an ideal of $F$ (i.e., of the ring of integers of 
$F$), let $I_\fm$ be the group of fractional ideals of $F$ prime to $\fm$, 
and let $k$ be a field of characteristic $0$. A homomorphism 
$\chi:I_\fm \to k^\times$ is an \emph{algebraic Hecke character} (of 
conductor $\leqslant \fm$) if there exists an algebraic homomorphism 
$\chi_\text{alg}:F^\times \to k^\times$ satisfying 
\begin{equation*}\tag{$*$}\label{VI:eq:5-3-1}
\text{For $x\in F^\times$, prime to $\fm$, totally positive, and 
$\equiv 1\pmod\fm$, one has $\chi((x))=\chi_\text{alg}(x)$.}
\end{equation*}

By the density of the set of $x$ in \eqref{VI:eq:5-3-1}, $\chi_\text{alg}$ 
is entirely determined by $\chi$. This is the \emph{algebraic part} 
of $\chi$. If $\chi((x))=\chi_\text{alg}(x)$ for $x$ totally positive and 
$\equiv 1\pmod{\fm'}$, $\chi$ extends to a Hecke character of 
conductor $\leqslant \inf(\fm,\fm')$: $I_{\fm+\fm'}\to k^\times$. We 
identify Hecke characters that coincide on their common domain of 
definition, and we call \emph{conductor} of $\chi$ the smallest $\fm'$ 
such that $\chi$ is of conductor $\leqslant \fm'$. 





\subsection{Remark}\label{VI:5-4}

If an algebraic Hecke character $\chi$ takes its values in a 
subfield $k'$ of $k$, it is already a Hecke character with values in 
$k'$: it suffices to observe that $\chi_\text{alg}:\res_{F/\dQ}(\dG_m) \to \dG_m$ 
is already defined over $k'$, and this follows from the Zariski density 
in $\res_{F/\dQ}(\dG_m)$ of the set of totally positive $x$ 
$\equiv 1\pmod\fm$. We can always take for $k'$ a subfield of $k$ of 
finite degree over $\dQ$. 





\subsection{}\label{VI:5-5}

If $\varepsilon$ is a totally positive unit $\equiv 1\pmod\fm$, we have 
$\chi_\text{alg}(\varepsilon) = \chi((\varepsilon)) = 1$. The homomorphism 
$\chi_\text{alg}$ therefore factors through the quotient $T_\fm$ of 
$\res_{F/\dQ}(\dG_m)$ by the Zariski closure of the group 
$E_\fm\subset F^\times$ of totally positive units $\equiv 1\pmod\fm$. 

According to Serre \cite[II.3]{se68}, if $k$ is an algebraic closure of 
$\dQ$, and $\fm$ is large enough, the characters $\prod \omega^{n_\omega}$ 
of $\res_{F/\dQ}(\dG_m)$ of the form $\chi_\text{alg}$ are characterized 
as follows: there must exist an integer $N$, the \emph{weight} of $\chi$ (or 
$\chi_\text{alg}$), such that for any element $\sigma$ of $\gal(k/\dQ)$ 
conjugate to complex conjugation, we have $n_\omega+n_{\sigma \omega}=N$. If 
$F_1$ is the largest subfield of $F$ which is a quadratic extension 
that is totally imaginary over a totally real field $F_1'$, this amounts to 
  saying that $\chi_\text{alg}$ is of the form $\chi_1\circ N_{F/F_1}$, and that 
  $\chi_1|F_1'=(N_{F_1'/\dQ})^N$. 

If $\chi$ is of weight $N$, for any ideal $\fa$ of $F$, $\chi(\fa)$ is a 
algebraic number all of whose conjugates have absolute value 
$N(\fa)^{N/2}$ \cite[II.3 prop.2]{se68}. 





\subsection{}\label{VI:5-6}

For a fixed number field $k$, additional conditions are 
imposed on $\chi_\text{alg}$. For any ideal $\fa$ of $F$ prime to the 
conductor, we have in fact 
\begin{equation*}\tag{5.6.1}\label{VI:eq:5-6-1}
  \chi_\text{alg}(\fa) = (\chi(\fa)) \text{,} 
\end{equation*}
so that $\chi_\text{alg}(\fa)$ is principal. Since the group of 
  ideal classes of $k$ is torsion, it is enough to prove the $n$-th power 
of \eqref{VI:eq:5-6-1}, for a suitable $n\ne 0$; this allows us 
to replace $\fa$ by $\fa^n$, therefore to assume $\fa=(x)$, with $x$ totally 
positive $\equiv 1\pmod\fm$. All that remains is to use the definitions. 

Together with \ref{VI:5-5}, this shows that $\chi_\text{alg}$ determines the 
valuations of $\chi(\fa)$ at all places of $k$. 





\subsection{}\label{VI:5-7}

Serre's group $S_\fm$ can be characterized as the 
group of multiplicative type whose group of characters (on any 
field) is the group of algebraic Hecke characters of conductor 
$\leqslant \fm$ (see \cite[II 2.1 et 2.2]{se68}). Its relation with 
$\ell$-adic representations is explained in \cite[II 2.3]{se68}.





\begin{theorem_}[{\cite{se68}}]\label{VI:5-8}
Let $\chi$ be an algebraic Hecke character of $F$ with values in $E_\lambda$, for 
$E_\lambda$ a finite extension of $\dQ_\ell$. There is then a unique 
  homomorphism $\chi_\lambda:\gal(\bar F/F)^\textnormal{ab} \to E_\lambda^\times$, such 
that 
\begin{enumerate}[\indent (i)]
  \item $\chi_\lambda$ is unramified outside the conductor $\ff$ of 
    $\chi$ and outside $\ell$. 
  \item For $\fp$ a prime ideal of $F$ prime to $\ff$ and to $\ell$, 
    and $F_\fp\in \gal(\bar F/F)^\textnormal{ab}$ the geometric Frobenius 
    at $\fp$, we have 
    \[
      \chi_\lambda(F_\fp) = \chi(\fp) \text{.}
    \] 
\end{enumerate}
\end{theorem_}





\subsection{}\label{VI:5-9}

At the end of this section, we give a criterion for a 
$\lambda$-adic representation to arise from a Hecke character, 
and in the next section we apply it to Jacobi sums. We 
thereby recover, in a somewhat generalized form, Weil's results 
\cite{we52,we74-2}, with the second part of Weil's proof 
replaced by a $\ell$-adic cohomology argument. 





\begin{theorem_}\label{VI:5-10}
Let $F$ and $k$ be two number fields, let $\lambda$ be a place of $k$ of 
residue characteristic $\ell$, and let 
$\chi_\lambda:\gal(\bar F/F)^\textnormal{ab}\to k_\lambda^\times$ be a 
homomorphism, unramified outside $\ell$ and a finite set $S$ of 
places. We assume that there exists a set $T$ of places (containing $S$ and the 
places above $\ell$), of density $0$, and an algebraic homomorphism 
$\chi_0:F^\times \to k^\times$ such that 
\begin{enumerate}[(i)]
  \item For $\bar k$ an algebraic closure of $k$, 
    $\chi_0:F^\times \to k^\times \hookrightarrow \bar k^\times$ is of the type 
    considered in \ref{VI:5-5}, of weight $N$; 
  \item For $\fp$ prime $\notin T$, $\chi_\lambda(F_\fp)$ lies in 
    $k^\times$, $(\chi_\lambda(F_\fp))=\chi_0(\fp)$, and all complex conjugates 
    of $\chi_\lambda(F_\fp)$ have absolute value $(N \fp)^{N/2}$. 
\end{enumerate}

Then $\chi_\lambda$ is defined (\ref{VI:5-8}) by an algebraic Hecke character 
of $F$ with values in $k$, whose algebraic part is $\chi_0$. 
\end{theorem_}

Let $\chi'$ be an algebraic Hecke character with values in a 
finite Galois extension $k'$ of $k$, of algebraic part $\chi_0$ 
(\ref{VI:5-5}). After enlarging $T$ if necessary, we may suppose that the conductor of 
$\chi'$ is supported in $T$. Let $\lambda'$ be a place of $k'$ above 
$\lambda$, and let $\chi_{\lambda'}$ be the composite 
$\gal(\bar F/F)^\text{ab}\to k_\lambda^\times \to {k'}_{\lambda'}^\times$, and 
$\chi_{\lambda'}':\gal(\bar F/F)^\text{ab}\to {k'}_{\lambda'}^\times$ defined 
by $\chi'$ (\ref{VI:5-8}). Set 
$\varepsilon = \chi_{\lambda'} \cdot {\chi'}_{\lambda'}^{-1}$.

Hypothesis (ii), together with \ref{VI:5-5}, \ref{VI:5-6}, and \ref{VI:5-8}, ensures that 
for $\fp\notin T$, we have $\varepsilon(F_\fp)\in {k'}^\times$, and that in every 
completion of $k'$ this number has absolute value $1$. It follows that 
$\varepsilon(F_\fp)$ belongs to the finite group $\dmu$ of roots of unity 
in $k'$. By the Chebotarev density theorem and the density-zero hypothesis 
on $T$, the $F_\fp$ ($\fp\notin T$) are dense in 
  $\gal(\bar F/F)^\text{ab}$, so that $\varepsilon(\sigma)\in \dmu$ for all 
$\sigma\in \gal(\bar F/F)^\text{ab}$: the character $\varepsilon$ is a 
  finite order character $\gal(\bar F/F)^\text{ab}\to \dmu$. By 
class field theory, it corresponds to a finite-order character 
of the idèle class group of $F$, i.e. to an algebraic Hecke character 
with trivial algebraic part. After correcting $\chi'$ by this 
character, we are reduced to the case in which 
$\chi_{\lambda'}=\chi_{\lambda'}'$, and it remains to show that $\chi'$, 
a priori taking values in $k'$, in fact takes values in $k$. If 
$\sigma\in \gal(k'/k)$, $\chi'$ and ${\chi'}^\sigma$ coincide on every prime 
ideal $\fp\notin T$, since 
$\chi'(\fp) = \chi'_{\lambda'}(F_\fp) = \chi_{\lambda'}(F_\fp) \in k^\times$. 
The characters $\chi_{\lambda'}'$, and ${\chi_{\lambda'}'}^\sigma$, coincide 
on a dense subset of $\gal(\bar F/F)^\text{ab}$ and hence are equal. Thus 
$\chi'={\chi'}^\sigma$, so $\chi'$ takes values in $k$. 










\section{Hecke characters defined by Jacobi sums}\label{VI:6}

The idea of this section is as follows. Consider a Gauss sum 
$\tau_{\dF_q}(\chi,\psi)$. Both this sum and the corresponding sheaf 
$\sF(\chi^{-1}\psi)$ are tied in two ways to the characteristic $p$ and to the 
base field $\dF_q$: 
\begin{enumerate}[\indent a)]
  \item because the additive character $\psi$ and the sheaf 
    $\sf(\psi)$ occur in it; 
  \item because $\sF(\chi)$ is defined from the exact Lang sequence 
    \[\xymatrix{
      0 \ar[r] 
        & k^\times \ar[r] 
        & V^\ast \ar[r] 
        & V^\ast \ar[r] 
        & 0 \text{.} 
    }\]
\end{enumerate}
These present two difficulties if we wish to lift the corresponding 
cohomological constructions to a number field $F$ and find 
$\ell$-adic representations of $\gal(\bar F/F)$ whose Frobenius eigenvalues 
are Gauss sums. 

If $\chi$ is trivial on $\dF_q^\times$, the sum $\tau_{\dF_q}(\chi,\psi)$ 
is independent of $\psi$ and is $q$ times a Jacobi sum. The latter is 
related to the cohomology of $W^\ast/\dG_m$, with coefficients in a sheaf 
obtained from $\chi$: thus $\psi$, and with it difficulty a), disappear. To 
resolve b), it suffices to describe the sheaves in question by means of 
exact K\"ummer sequences (which make sense in every characteristic). This is 
possible by \ref{VI:4-9}. 

Once the $\ell$-adic representations have been found, we use 
\ref{VI:5-10} and Stickelberger's theorem to show that they 
come from algebraic Hecke characters. 





\subsection{}\label{VI:6-1}

If $E$ is a field with algebraic closure $\bar E$, and $\lambda$ is 
a character, of order dividing $n$, of $\mu_d(\bar E)$, we denote by $_n\lambda$ the 
character of $\mu_n(\bar E)$ such that $\lambda$ and $_n \lambda$ induce the 
same character of $\widehat\dZ(1)_{\bar E} = \varinjlim \mu_n(\bar E)$. The 
same notation is used when $\lambda$ is a character of $\widehat\dZ(1)_{\bar E}$.

Let $F$ and $k$ be two number fields, $\bar F$ an algebraic closure of 
$F$, $I$ a finite set equipped with an action of $\gal(\bar F/F)$ and 
$\lambda=(\lambda_i)_{i\in I}$ a family of characters 
$\lambda_i:\widehat\dZ(1)_{\bar F} \to k^\times$. It is assumed that 
$\lambda_i\ne 1$ ($i\in I$), that the product of $\lambda_i$ is trivial and that 
the family $\lambda$ is defined over $F$, i.e. that for every 
$\sigma\in \gal(\bar F/F)$, we have 
$\lambda_{\sigma(i)}=\lambda_i\circ \sigma^{-1}$. 

The Galois set $I$ is the set of homomorphisms into $\bar F$ of an 
algebra $E$ separable over $F$. If $C$ is the set of orbits of 
$\gal(\bar F/F)$ in $I$, we have a decomposition of $E$ as a product of fields 
$E=\prod_{\alpha\in C} E_\alpha$, where $\alpha$ is identified with the set of 
embeddings of $E_\alpha$ in $\bar F$. If the characters $\lambda_i$, for 
$i\in \alpha$, are of order $d_\alpha$, the field $E_\alpha$ contains the 
$d_\alpha$-th roots of $1$ and there exists 
$\lambda_\alpha:\mu_{d_\alpha}(E_\alpha) \hookrightarrow k^\times$ such that, if 
$i\in \alpha$ corresponds to the embedding $\sigma_i$ of $E_\alpha$ in $\bar F$, 
then $_{d_\alpha} \lambda_i = \lambda_\alpha \sigma_i^{-1}$. 





\subsection{}\label{VI:6-2}

Let $n>0$ be a multiple of the $d_\alpha$, let $\Sigma$ be the set of places of $F$ 
at which $E/F$ ramifies or which divide $n$, let $\mu$ be a place of $k$ of 
residue characteristic $\ell$, and let $A$ be the ring of elements of $F$ 
integral outside $\Sigma$ and $\ell$. The integral closure $M$ of $A$ 
is finite étale of rank $N=|I|$ over $A$. With the notation of \ref{VI:4-6}, 
\ref{VI:4-8}, $_n \lambda_i$ define 
\[
  _n\lambda:(V^\ast/\dG_m)_n \to k^\times \subset k_\mu^\times \text{.} 
\]

Let $a$ be the projection of $W^\ast/\dG_m$ onto $\spec(A)$. By 
\ref{VI:4-20}, the sheaf $\eR^{N-2} a_! \sK_n(_n \lambda)$ is smooth of rank 
one; it defines an $\ell$-adic representation 
\[
  j_\mu[\lambda]:\gal(\bar F/F)^\text{ab} \to k_\mu^\times 
\]
  (or simply $j_\mu$) unramified outside $\Sigma$ and $\ell$. 





\subsection{}\label{VI:6-3}

Let us compute $j_\mu(F_\fp)$. By \ref{VI:4-20}, it suffices to do so 
after reduction modulo $\fp$. Put $f=A/\fp$, $e=M/\fp  M$, and let 
$\bar f$ be the algebraic closure of $f$ determined by a place of $\bar F$ 
above $\fp$. We still have $I=\hom_f(e,\bar f)$. After reduction, 
$_n\lambda$ still has a description as in \ref{VI:6-1}: 
\begin{enumerate}[a)]
  \item Let $D$ be the set of orbits of $\gal(\bar f/f)$ in $I$. The 
    decomposition of $e$ as a product of fields has the form 
    $e=\prod_{\beta\in D} e_\beta$, where $\sigma_i$ ($i\in \beta$) 
    are identified with the embeddings of $e_\beta$ into $\bar f$. For $\beta\in D$, 
    denote by $\alpha(\beta)$ the element of $C$ containing $\beta$, and put 
    $d_\beta=d_{\alpha(\beta)}$. Let 
    \[\xymatrix{
      \lambda_\beta : \mu_{d_\beta}(e_\beta) 
        & \ar[l]_-\sim \mu_{d_\beta}(E_{\alpha(\beta)}) \ar[r]^-{\lambda_{\alpha(\beta)}} 
        & k^\times \text{.} 
    }\]
  \item If $i\in \beta$ corresponds to the embedding $\sigma_i$ of $e_\beta$ into 
    $\bar f$, the reduction 
    $_{d_\beta} \bar\lambda_i:\mu_{d_\beta}(\bar f) \to k^\times$ 
    of $_{d_\beta} \lambda_i$ is 
    $\lambda_\beta \circ \sigma_i^{-1}$. 
\end{enumerate}

Let $q_\beta$ be the number of elements of $e_\beta$, 
  $\chi_\beta={}_{q_\beta-1}\lambda_\beta$, and 
$\chi:e^\times \to k^\times$ the character with coordinates $\chi_\beta$. On the 
affine space over $A/\fp$ defined by $e$ (\ref{VI:4-6}), or rather on the torus 
$e_\beta^\times$, we have isomorphisms 
  $\sK_n((_n\bar\lambda_i)_{i\in\beta}) = \sK_{d_\beta}((\lambda_\beta \circ \sigma_i^{-1})_{i\in\beta}) =\sK_{q_\beta-1}((\chi_\beta\circ \sigma_i^{-1})_{i\in\beta}) = \sF(\chi_\beta)$ (\ref{VI:4-9}). 
  On $e^\times$, we therefore have $\sK_n((_n\bar\lambda_i)_{i\in I}) = \sF(\chi)$, 
and this isomorphism descends to the reduction modulo $\fp$ of $V^\ast/\dG_m$. 
It follows that





\begin{proposition_}\label{VI:6-4}
With the notations of \ref{VI:6-2} and \ref{VI:6-3}, $j_\mu(F_\fp)$ is the 
Jacobi sum $J(\chi)$. 
\end{proposition_}

Let us rewrite the formulas linking $(\lambda_i)_{i\in I}$ to 
  $(\chi_\beta)_{\beta\in D}$: for $i\in \beta\subset \alpha$, 
\begin{align*}
  _{d_\alpha} \lambda_i &= \lambda_\alpha\circ \sigma_i^{-1} && \lambda_\alpha:\mu_{d_\alpha}(E_\alpha) \to k^\times \\
  \lambda_\beta &= \text{reduction of $\lambda_\alpha$} && :\mu_{d_\alpha}(e_\beta) \to k^\times \\
  \chi_\beta &= {}_{q_\beta-1}\lambda_\beta: x\mapsto \lambda_\beta(x^{(q_\beta-1)/d_\alpha}) && :e_\beta^\times \to k^\times 
\end{align*}





\begin{theorem_}\label{VI:6-5}
The representation $j_\mu$ is defined by an algebraic Hecke character 
of $F$ with values in $k^\times$. 
\end{theorem_}

We shall verify this using \ref{VI:5-10}. Here we follow Weil's 
proof \cite{we52,we74-2}, using Stickelberger's theorem, 
and I shall give only a few indications. 





\subsection{}\label{VI:6-6}

Given $j:\gal(\bar F/F) \to k_\lambda^\times$ and a finite extension 
$F'/F$, one verifies using \ref{VI:5-10} that $j$ is defined by a 
Hecke character if and only if $j|\gal(\bar F/F')$ is. This 
remark allows us, in \ref{VI:6-5}, to reduce to the case where Galois acts 
trivially on $I$ and where $F$ contains the $n$-th roots of 
unity. We may then reduce to the case $F=\dQ(\sqrt[n]{1})$ and 
take $k=F$. In this case, each $_n \lambda_i$ is defined by 
$a_i\in \dZ/n$: it is $x\mapsto x^{a_i}:\mu_n \to k^\times$. We have 
$\sum a_i=0$, and $a_i\ne 0$. 





\subsection{}\label{VI:6-7}

In a number field, the places of degree greater than $1$ always form a 
set of density $0$. By applying \ref{VI:5-10}, we may therefore disregard 
them and use Stickelberger only for a prime field. Let 
$p$ be prime, let $0<a<p-1$, and let $\widetilde x:\dF_p^\times \to \dQ_p^\times$ be the 
``multiplicative lift'' character; put 
  $g=-\sum_{x\in \dF_p^\times} \widetilde x^{-a} \zeta^x\in \dQ_p(\zeta)$, for 
$\zeta$ a $p$-th root of $1$. Again denote by $x$ the integer in $[0,p-1]$ 
  lifting $x$, and expand $\zeta^x=(1+\pi)^x$ by the binomial formula. 
  We find $g=-\sum_{i=0}^{p-1} A_i\pi^i$ with $A_i\in \dZ_p$, whose reduction 
$\mod p$ is the sum $\sum_{x\in \dF_p^\times} x^{-a} \binom x i$. For $b$ not 
divisible by $p-1$, we have $\sum_{x\in \dF_p^\times} x^b=0$. If $i<a$, 
$\binom x i$ is a polynomial in $x$ of degree $i<a$, whence 
$A_i\equiv 0\pmod p$. Likewise, $A_a\equiv \frac{p-1}{a!} \pmod p$, and 
$g\sim \frac{\pi^a}{a!}$: the valuation of $g$ is given by 
$v(g)/v(p) = \frac{a}{p-1}$. 





\subsection{}\label{VI:6-8}

Let $\sigma_i$ ($i\in (\dZ/n)^\times$) be the element of 
$\gal(\dQ(\sqrt[n]{1})/\dQ)$ such that $\sigma_i(\zeta)=\zeta^i$ for $\zeta^n=1$, 
and write the group of algebraic homomorphisms 
$\dQ(\sqrt[n]{1})^\times \to \dQ(\sqrt[n]{1})^\times$ additively. For $\lambda$ as in 
\ref{VI:6-6}, defined by a family 
$(a_i)_{i\in I}$ of nonzero integers $\mod n$ whose sum is zero, 
\ref{VI:5-10}, \eqref{VI:eq:4-15-1} and \ref{VI:6-7} 
show that $j_\mu[\lambda]$ is defined by an algebraic Hecke character 
$j[a]$ and that, if $N$ denotes the norm, its algebraic part is 
given by 
\begin{equation*}\tag{6.8.1}\label{VI:eq:6-8-1}
  (N j[a])_\text{alg} = \sum_{i\in (\dZ/n)^\times} \sum_{j\in I} \left\langle \frac{i a_j}{n} \right\rangle \sigma_i^{-1} 
\end{equation*}
  where $\langle \cdot \rangle$ is the fractional part. 





\subsection{}\label{VI:6-9}

The most interesting case is that in which $F\subset \dQ(\sqrt[n]{1})$ and 
the family of $_n \lambda_i$ is constructed as follows: 
\begin{enumerate}[a)]
  \item Take a finite family $(A_j)_{j\in J}$ of orbits of 
    $\gal(\dQ(\sqrt[n]{1})/F)=H\subset (\dZ/n)^\times$ in $\dZ/n$. Assume that 
    $A_j\ne\{0\}$ and that $\sum_{j\in J} \sum_{a\in A_j} a=0$. 
  \item Take $I=\sqcup_{j\in J} A_j$; Galois acts on each $A_j$. 
  \item Take $\bar F\supset \dQ(\sqrt[n]{1})$, $k=\dQ(\sqrt[n]{1})$ and, if 
    $i\in I$ corresponds to $a\in A_j\subset (\dZ/n)^\times$, 
    $_n\lambda_i(x)=x^a$. 
\end{enumerate}

Using \ref{VI:4-20}, one verifies that a general $j_\mu$ (relative to 
$F_1$) is obtained from a $j_\mu$ as above (for 
$F\subset \dQ(\sqrt[n]{1})$, with $F_1$ an extension of $F$) by restriction to 
$\gal(\bar F/F_1)$ and multiplication by a character of order two. 

For $a\in \gal(k/\dQ)=(\dZ/n)^\times$, $a$ transforms the family of $A_j$ 
into the family of $a A_j$. If $a\in H$, the family of $A_j$ is therefore preserved, 
and the $j_\mu(F_\fp)$ are invariant under $H$: 





\begin{proposition_}\label{VI:6-10}
Under the assumptions of \ref{VI:6-9}, $j_\mu$ is defined by a character 
that is an algebraic Hecke character of $F$ with values in $F$. 
\end{proposition_}

Let us again take $F=\dQ(\sqrt[n]{1})$ and consider the Hecke characters 
$j[a]$ defined by a family $(a_i)_{i\in I}$, $i\in \dZ/n$, 
$a_i\ne 0$, $\sum a_i=0$. The following theorem was obtained 
independently by Vishik. 





\begin{theorem_}\label{VI:6-11}
Every algebraic Hecke character of $\dQ(\sqrt[n]{1})$ has a power 
in the group generated by the $j[a]$ and the norm. 
\end{theorem_}

Since we work ``up to torsion,'' only the algebraic part of the 
Hecke characters under consideration matters. Applying \eqref{VI:eq:6-8-1} to 
$j[a]$, with $a=$ an element of $\dZ/n$ repeated $n$ times, we 
are reduced to the following lemma.





\begin{lemma_}\label{VI:6-12}
Write $\sum x_j \sigma_j$ for the elements of the group algebra 
  $\dQ[(\dZ/n)^\times]$, and let $X$ be the subgroup consisting of those for which 
$x_j+x_{-j}=C^{\underline{t e}}$. Then $X$ is generated over $\dQ$ by 
$N=\sum \sigma_j$ and by the elements 
  $g_a=\sum \left\langle \frac{i a}{n}\right\rangle \sigma_i^{-1}$ 
($a\in \dZ/n$, $a\ne 0$). 
\end{lemma_}

Both $X$ and $Y=\langle N,(g_a)\rangle$ are ideals of the group 
algebra, and it suffices to check that they are annihilated by the same 
complex characters $\chi$ of $(\dZ/n)^\times$. This amounts to saying that a 
odd character $(\chi(j)=-\chi(-j)$) never annihilates all the $g_a$. If 
$\chi$ is primitive, we have 
\[
  \chi(g_1) = \sum \frac i n \chi^{-1}(i) = -L(\chi,0)\ne 0 \text{.} 
\]
In the general case, if $\chi$ comes from a primitive character of 
  $(\dZ/d)^\times$, $\chi(g_{n/d})$ reduces to an analogous sum over 
  $(\dZ/d)^\times$, hence the lemma. 










\section{Generalized Kloosterman sums}\label{VI:7}





\subsection{}\label{VI:7-1}

Let $\dF_q$ be a finite field with $q$ elements, and let $\psi$ be a nontrivial 
additive character. In this section, we undertake a cohomological study of 
generalized Kloosterman sums 
\begin{equation*}\tag{7.1.1}\label{VI:eq:7-1-1}
  \kloos_{n,a} = \sum_{x_1\dotsm x_n = a} \psi(x_1+\cdots + x_n) \qquad (n\geqslant 1) \text{.} 
\end{equation*}
For $a=0$, this sum is elementary (see \ref{VI:7-7}): 
\begin{equation*}\tag{7.1.2}\label{VI:eq:7-1-2}
  \kloos_{n,0} = (-1)^{n-1} \text{.} 
\end{equation*}
The interesting case is $a\ne 0$. The $x_i$ then belong to 
$\dF_q^\times$. In this case we shall prove the bound 
\begin{equation*}\tag{7.1.3}\label{VI:eq:7-1-3}
  \left|\kloos_{n,a}\right| \leqslant n q^{\frac{n-1}{2}} \text{.} 
\end{equation*}

Our essential tools will be the cohomological counterparts of the 
following identities: 
\begin{align*}\tag{7.1.4}\label{VI:eq:7-1-4}
  \kloos_{n,a} &= \sum_{x_1} \psi(x_1) \sum_{x_2\dotsm x_n = a/x_1} \psi(x_2+\cdots + x_n) \\ \notag
    &= \sum_{x\in \dF_q^\times} \psi(x) \kloos_{n-1,a/x} \qquad (a\ne 0,n\geqslant 2) \\ \tag{7.1.5}\label{VI:eq:7-1-5}
  \sum_a \kloos_{n,a} &= \sum_a \sum_{x_1 \dotsm x_n = a} \psi\left(\sum x_i\right) = \sum_{x_1,\dots,x_n} \psi\left(\sum x_i\right) = 0 
\end{align*}
For $\chi$ a nontrivial character of $\dF_q^\times$, 
\begin{equation*}\tag{7.1.6}\label{VI:eq:7-1-6}
  \sum \chi(a) \kloos_{n,a} = \sum_{x_1,\dots,x_n} \chi\left(\prod x_i\right) \cdot \psi\left(\sum x_i\right) = (-\tau(\chi^{-1},\psi))^n 
\end{equation*}
(Gauss sum). 





\subsection{}\label{VI:7-2}

Let $k$ be an étale algebra of degree $n$ over $\dF_q$, and let $\varepsilon(k)$ 
be as in \eqref{VI:eq:4-5-1}. Put 
\[
  \kloos_{k,a} = \sum_{N_{k/\dF_q}(x)=a} \psi \tr(x) \text{.} 
\]
We likewise have 
\begin{align*}\tag{7.2.1}\label{VI:eq:7-2-1}
  \kloos_{k,0} &= (-1)^{n-1} \cdot \varepsilon(k) \\ \tag{7.2.2}\label{VI:eq:7-2-2}
  \sum_a \kloos_{k,a} &= 0 \\ \tag{7.2.3}\label{VI:eq:7-2-3}
  \sum_a \chi(a) \kloos_{k,a} &= (-1)^n \tau_{\dF_q}\left(\chi^{-1}\circ N_{k/\dF_q},\psi\right) \text{.} 
\end{align*}

By Fourier inversion on $\dF_q^\times$, we deduce an expression of 
$\kloos_{k,a}$ in terms of Gauss sums: 
\begin{align*} 
  \kloos_{k,a} &= \frac{1}{q-1} \sum_\chi \chi(a) \sum_{b\in \dF_q^\times} \chi^{-1}(b) \kloos_{k,b} \\ \tag{7.2.4}\label{VI:eq:7-2-4}
  (-1)^n \kloos_{k,a} &= \frac{1}{q-1} \left(\varepsilon(k) + \sum_{\chi\ne 1} \chi(a) \tau_{\dF_q}\left(\chi\circ N_{k/\dF_q},\psi\right) \right) \text{.} 
\end{align*}

The Hasse--Davenport identity now makes it possible to compare $\kloos_{k,a}$ 
with $\kloos_{n,a} = \kloos_{\dF_{q^n},a}$: 
\begin{equation*}\tag{7.2.5}\label{VI:eq:7-2-5}
  \kloos_{k,a} = \varepsilon(k) \kloos_{n,a} \text{.} 
\end{equation*}

I know of no cohomological interpretation of \eqref{VI:eq:7-2-4}, but 
I shall give one for \eqref{VI:eq:7-2-5}. Let us return to the sums $\kloos_{n,a}$. 





\subsection{}\label{VI:7-3}

Let $\dF$ be an algebraic closure of $\dF_q$ and, for $a\in \dF$, let 
$V_a^{n-1}$, or simply $V_a$, be the hypersurface in $\dA^n$ with equation 
$x_1 \dotsm x_n=a$. As usual, we shall regard $\psi$ as taking 
$\ell$-adic rather than complex values. With the notation of 
\ref{VI:1-8}(ii), \eqref{VI:eq:7-1-3} follows from 





\begin{theorem_}\label{VI:7-4}
The cohomology of $V_a^{n-1}$ with coefficients in $\sF(\psi(\sum x_i))$ satisfies 
\begin{enumerate}[\indent (i)]
  \item $\h_c^i=0$ for $i\ne n-1$; 
  \item $\h_c^\bullet{} \iso \h^\bullet{}$; 
  \item for $a\ne 0$, $\dim \h_c^{n-1}{}=n$; 
  \item for $a=0$, $\h_c^{n-1}{}$ is canonically isomorphic to $E_\lambda$. 
\end{enumerate}
\end{theorem_}





\subsection{}\label{VI:7-5}

As usual, this theorem also controls the dependence of 
$\kloos_{n,a}$ on $q$: for each $a\in \dF_q^\times$, there are $n$ Frobenius 
eigenvalues $\alpha_1,\dots,\alpha_n$, of complex absolute value 
$q^{\frac{n-1}{2}}$, such that 
\[
  \sum_{\substack{x_1\dotsm x_n = a \\ x_i\in \dF_{q^m}}} \psi\circ \tr\left(\sum x_i\right) = (-1)^{n-1} \left(\alpha_1^m + \cdots + \alpha_n^m \right) \text{.} 
\]

For even $n$, $x\mapsto -x$ is an involution of $V_a$, which transforms 
$\sF(\psi)$ into its dual. Arguing as in \ref{VI:3-6}, we deduce a 
canonical nondegenerate alternating form on $\h_c^{n-1}(V_a,\sF(\psi))$, 
with values in $\dQ_\ell(-(n-1))$, and hence the following. 





\begin{corollary_}\label{VI:7-6}
With the notations of \ref{VI:7-5}, if $n$ is even, the $\alpha_i$ group together 
into $\frac n 2$ pairs of roots $\alpha,\bar\alpha$ whose product is 
$q^{n-1}$. 
\end{corollary_}





\subsection{}\label{VI:7-7}

Let us verify the case $a=0$ of \ref{VI:7-5}. The hypersurface $V_0$ is the union 
of the coordinate hyperplanes $x_i=0$. Let us compute the Leray spectral sequence 
\eqref{VI:eq:2-6-1*} for this covering of $V_0$, for cohomology with or 
without supports and with coefficients in $\sF(\psi)$. By \ref{VI:2-7*}*, all 
the initial terms vanish except the cohomology of the intersection $\{0\}$ of 
all the coordinate hyperplanes. What remains is an isomorphism 
$\h_c^0{}\iso \h^0{}=E_\lambda$ which justifies \ref{VI:7-4}. 

We shall prove \ref{VI:7-4} and \ref{VI:7-8} below by simultaneous induction 
on $n$, using the spectral sequence of which 
\eqref{VI:eq:7-1-4} is the counterpart. This requires control, as $a$ varies, 
of $\kloos_{n,a}$. 

Let $\pi:\dA^n \to \dA$ be the map ``product of the coordinates,'' and let 
$\sigma$ be the sum of the coordinates. The hypersurfaces $V_a$ are the fibers of 
$\pi$.





\begin{theorem_}\label{VI:7-8}
\begin{enumerate}[(i)]
  \item The sheaf $\eR^{n-1}\pi_! \sF(\psi\sigma)$ is smooth of rank $n$ on 
    $\dA^1\setminus \{0\}$. 
  \item Its extension by $0$ to $\dP^1$ is the direct image of its 
    restriction to $\dA^1\setminus \{0\}$. 
  \item At $0$, the monodromy is unipotent, with a single Jordan block.
  \item At $\infty$, wild inertia acts without fixed points $\ne 0$, and the 
    Swan conductor is $1$. 
  \item We have $\eR^i \pi_! \sF(\psi\sigma) \iso \eR^i \pi_\ast \sF(\psi\sigma)$ 
    (zero for $i\ne n-1$).
\end{enumerate}
\end{theorem_}
\begin{proof}[Proof]
that \ref{VI:7-4} (for a given value of $n$) implies \ref{VI:7-8} 
(for the same value of $n$). 

The fiber of $\eR^i \pi_! \sF(\psi\sigma)$ at $a$ is 
$\h_c^i(V_a,\sF(\psi\sigma))$, and on the dense open subset of $\dA^1$ where the 
  $\eR^i \pi_!\sF(\psi\sigma)$ are locally constant, the fiber of 
$\eR^i \pi_\ast \sF(\psi\sigma)$ is $\h^i(V_a,\sF(\psi\sigma))$ 
(\hyperref[VII]{Finiteness\ Theorem} \ref{VII:2-1}). The hypothesis \ref{VI:7-4}($n$) 
therefore gives us the following. 





\begin{lemma_}\label{VI:7-9}
\begin{enumerate}[(i)]
  \item $\eR^i \pi_! \sF(\psi\sigma)=0$ for $i\ne n-1$. 
  \item For $i=n-1$, the fiber of this sheaf at any point $a\ne 0$ has 
    constant rank $n$. At $0$, it has rank $1$. 
  \item On a dense open $U$, we have 
    $\eR^i \pi_! \sF(\psi\sigma) \iso \eR^i \pi_\ast \sF(\psi\sigma)$. 
\end{enumerate}
\end{lemma_}





\subsection{}\label{VI:7-10}

We shall now use the Leray spectral sequence of $\pi$ for 
cohomology with or without compact supports, with coefficients in $\sF(\psi\sigma)$. 
The abutment is known: 
$\h_c^\bullet(\dA^n,\sF(\psi\sigma))=\h^\bullet(\dA^n,\sF(\psi\sigma)) = 0$ 
(\ref{VI:2-7*}*). For cohomology with compact supports, \ref{VI:7-9}(i) ensures that 
$E_2^{pq}=0$ for $q\ne n-1$. The terms $E_2$ are therefore all zero: 
\begin{equation*}\tag{7.10.1}\label{VI:eq:7-10-1}
  \h_c^p(\dA^1,\eR^{n-1} \pi_!\sF(\psi\sigma)) = 0 \text{.} 
\end{equation*}

  For $p=0$, this means that $\eR^{n-1}\pi_! \sF(\psi\sigma)$ has no sections 
with point support. In view of the constancy of rank in \ref{VI:7-9}(ii), we obtain 
\begin{equation*}\tag{7.10.2}\label{VI:eq:7-10-2}
  \begin{array}{c}
    \text{The sheaf $\eR^{n-1} \pi_! \sF(\psi\sigma)$ on $\dA^1$ is smooth on $\dA^1\setminus \{0\}$, and is } \\ 
    \text{a subsheaf of the direct image $\sG$ of its restriction to $\dA^1\setminus \{0\}$.}
  \end{array}
\end{equation*}

The long exact cohomology sequence of 
\[\xymatrix{
  0 \ar[r] 
    & \eR^{n-1} \pi_! \sF(\psi\sigma) \ar[r] 
    & \sG \ar[r] 
    & \sQ \ar[r] 
    & 0 \text{.} 
}\]
where the support of $\sQ$ is concentrated at $0$, gives 
\[\xymatrix{
  \h_c^0(\dA^1,\sG) \ar[r] 
    & \sQ_0 \ar[r] 
    & \h_c^1(\dA^1,\eR^{n-1} \pi_! \sF(\psi\sigma)) \text{.} 
}\]
Since the outer terms vanish, $\sQ=0$ and 
\begin{equation*}\tag{7.10.3}\label{VI:eq:7-10-3}
  \text{$\eR^{n-1} \pi_! \sF(\psi\sigma)$ is the direct image of its restriction to $\dA^1\setminus \{0\}$.}
\end{equation*}

Let $j$ be the inclusion of $\dA^1$ into $\dP^1$. The ``derived-category 
version'' of the Leray spectral sequences is 
\begin{align*}
  \hh^\bullet(\dP^1,j_! \eR\pi_! \sF(\psi\sigma)) &= \hh_c^\bullet(\dA^1,\eR \pi_! \sF(\psi\sigma)) = \h_c^\bullet(\dA^n,\sF(\psi\sigma)) = 0 \\
  \hh^\bullet(\dP^1,\eR j_\ast \eR \pi_\ast \sF(\psi\sigma)) &= \hh^\bullet(\dA^1,\eR \pi_\ast \sF(\psi\sigma)) = \h^\bullet(\dA^n,\sF(\psi\sigma)) = 0 \text{.}
\end{align*}
Let $\Delta$ be the mapping cylinder of 
$j_! \eR\pi_! \sF(\psi\sigma) \to \eR j_\ast \eR \pi_\ast \sF(\psi\sigma)$. 
According to \ref{VI:7-9}(iii), the cohomology sheaves of this complex of 
sheaves have finite support. On the other hand, the long exact sequence in 
cohomology of the triangle 
$(j_!\eR \pi_! \sF(\psi\sigma),\eR j_\ast \eR \pi_\ast \sF(\psi\sigma),\Delta)$ 
shows that $\hh^\bullet(\dP^1,\Delta)=0$. We have 
$\hh^\bullet(\dP^1,\Delta) = \h^0(\dP^1,\sH^\bullet(\Delta))$, and finally 
$\Delta=0$: 
\[\xymatrix{
  j_! \eR\pi_! \sF(\psi\sigma) \ar[r]^-\sim 
    & \eR j_\ast \eR \pi_\ast \sF(\psi\sigma) \text{.} 
}\]

In particular, $\eR\pi_! \sF(\psi\sigma)\iso \eR\pi_\ast \sF(\psi\sigma)$ and, 
since these complexes of sheaves have only one nonzero cohomology sheaf, 
\begin{equation*}\tag{7.10.4}\label{VI:eq:7-10-4}
\xymatrix{
  j_! \eR \pi_! \sF(\psi\sigma) \ar[r]^-\sim 
    & j_\ast \eR \pi_! \sF(\psi\sigma) \text{.} 
}
\end{equation*}
This completes the proof of (i)(ii)(v). 





\subsection{}\label{VI:7-11}

Let $\swan_0$ and $\swan_\infty$ be the Swan conductors at $0$ and at 
$\infty$ of $\eR^{n-1} \pi_! \sF(\psi\sigma)$. We have 
$\chi_c(\eR^{n-1} \pi_! \sF(\psi\sigma))=0$ \eqref{VI:eq:7-10-1} and the formula 
of Euler--Poincaré \eqref{VI:eq:3-2-1} reduces to 
\begin{equation*}\tag{7.11.1}\label{VI:eq:7-11-1}
  \swan_0 + \swan_\infty =1 \text{:} 
\end{equation*}
one of these conductors is $0$ (tame ramification), and the other is $1$. 





\subsection{}\label{VI:7-12}

Let $\sG$ be a sheaf of rank $1$ on $\dG_m$, of Kummer type 
\ref{VI:4-7} ($\sG=\sK_n(\chi)$) and nonconstant ($\chi\ne 1$). The sheaf 
$\sF(\psi\sigma)\otimes \pi^\ast \sG$ on 
$V^\ast=\pi^{-1}(\dG_m) = \dA^n\setminus \pi^{-1}(0)$ is then of the type 
encountered in the study of Gauss sums, and \ref{VI:4-11}(ii) gives us 
\begin{equation*}\tag{7.12.1}\label{VI:eq:7-12-1}
\xymatrix{
  \h_c^\bullet(V^\ast,\pi^\ast\sG\otimes \sF(\psi\sigma)) \ar[r]^-\sim 
    & \h^\bullet(V^\ast,\pi^\ast \sG\otimes \sF(\psi\sigma)) \text{.} 
}
\end{equation*}
The Leray spectral sequence of $\pi$ transforms this statement into 
\begin{equation*}\tag{7.12.2}\label{VI:eq:7-12-2}
\xymatrix{
  \h_c^\bullet(\dG_m,\sG\otimes \eR^{n-1} \pi_! \sF(\psi\sigma)) \ar[r]^-\sim 
    & \h^\bullet(\dG_m,\sG\otimes \eR^{n-1} \pi_\ast \sF(\psi\sigma)) \text{.} 
}
\end{equation*}

Let $i$ be the inclusion of $\dG_m$ into $\dP^1$, and let $\Delta$ be the mapping cylinder 
of 
  $\delta:i_!(\sG\otimes \eR^{n-1}\pi_! \sF(\psi\sigma))\to \eR i_\ast (\sG\otimes \eR^{n-1} \pi_! \sF(\psi\sigma))$
(where $\eR^{n-1} \pi_! = \eR^{n-1}\pi_\ast$). Reasoning as in 
\ref{VI:7-10}, we deduce from \eqref{VI:eq:7-12-2} that $\Delta=0$. In 
particular 
\begin{equation*}\tag{7.12.3}\label{VI:eq:7-12-3}
  \text{At $0$ and at $\infty$, inertia acts without fixed points $\ne 0$ on $\sG\otimes \eR^{n-1}\pi_! \sF(\psi\sigma)$.}
\end{equation*}

At $\infty$, this holds even if $\sG$ is trivial \eqref{VI:eq:7-10-4}, 
so wild inertia acts without a fixed point. In particular, the 
ramification is wild: by \ref{VI:7-11}, we have $\swan_\infty=1$, and the 
ramification at $0$ is tame. 

At $0$, \eqref{VI:eq:7-12-3} forces the ramification -- which is tame -- of 
$\eR^{n-1} \pi_! \sF(\psi\sigma)$ to be unipotent. By 
\eqref{VI:eq:7-10-3} and \ref{VI:7-9}(ii), there is a single Jordan block. This 
completes the proof of \ref{VI:7-8}($n$). 





\subsection{}\label{VI:7-13}

For $n=1$, \ref{VI:7-4} is trivial. To complete the proof of \ref{VI:7-4} and 
\ref{VI:7-8}, we still have to prove that \ref{VI:7-8}, for a given value 
of $n$, implies \ref{VI:7-4} for $n+1$. The case $a=0$ having been 
  treated (\ref{VI:7-7}), we assume that $a\ne 0$. We shall write $x_0,\dots,x_n$ for 
the coordinates of $\dA^{n+1}$. Let $g=x_0:V_a\to \dG_m$, $\tau$ 
the involution $x\mapsto a x^{-1}$ of $\dG_m$, and let $\pi$, as above, be the 
product of the coordinates: $\dA^{n+1} \to \dA^1$, with $\sigma$ their sum. We 
shall denote by $\sF(\psi\sigma)$ the sheaf $\sF(\psi(\sum x_i))$ both on $\dA^{n+1}$ 
and on $\dA^n$. 





\subsection{}\label{VI:7-14}

Let us write down the Leray spectral sequence for $g$. Since the sheaf $\sF(\psi\sigma)$ on 
$\dA^{n+1}$ is the external tensor product of the sheaves $\sF(\psi)$ on 
$\dA^1$ and $\sF(\psi\sigma)$ on $\dA^n$, we find (cf. \eqref{VI:eq:7-1-4}). 
\begin{align*}
  'E_2^{p q} &= \h_c^p(\dG_m,\sF(\psi)\otimes \tau^\ast \eR^q \pi_! \sF(\psi\sigma)) \Rightarrow \h_c^{p+q}(V_a,\sF(\psi\sigma)) \\
  ''E_2^{p q} &= \h^p(\dG_m,\sF(\psi)\otimes \tau^\ast \eR^q \pi_\ast \sF(\psi\sigma)) \Rightarrow \h^{p+q}(V_a,\sF(\psi\sigma)) \text{.} 
\end{align*}
By hypothesis, 
$\eR^q \pi_! \sF(\psi\sigma) \iso \eR^q \pi_\ast \sF(\psi\sigma)$, and this 
sheaf is zero for $q\ne n-1$. The sheaf $\sF(\psi)$ on $\dG_m$ is 
wildly ramified at $\infty$ and unramified at $0$, while 
$\tau^\ast \eR^q \pi_! \sF(\psi\sigma)$ is wildly ramified at $0$ (without 
invariants under wild inertia, and with Swan conductor $1$) and tamely 
ramified at $\infty$. Their tensor product is therefore 
\begin{enumerate}[\indent a)]
  \item wildly ramified at $0$ and $\infty$, without invariants under wild 
    inertia; 
  \item of Swan conductor $1$ at $0$ and $n$ (the rank of 
    $\eR^q \pi_! \sF(\psi\sigma))$ at $\infty$. 
\end{enumerate}
We have $'E_2^{p q} \iso ''E_2^{p q}$, and $'E_2^{p q}=0$ except for $p=1$ and 
$q=n-1$. The Euler-Poincaré formula finally gives 
\[
  \dim{'E_2^{1,n-1}} = 1+n \text{;} 
\]
this completes the proof. 
\end{proof}





\subsection{Remark}\label{VI:7-15}

The isomorphism obtained in \ref{VI:7-14} 
\begin{equation*}\tag{7.15.1}\label{VI:eq:7-15-1}
  \h_c^n(V_a^n,\sF(\psi\sigma)) = \h_c^1(\dG_m,\sF(\psi)\otimes \tau^\ast \eR^{n-1} \pi_! \sF(\psi\sigma)) 
\end{equation*}
allows us, with the notation of \ref{VI:7-5}, to calculate the product $\det(F)$ 
of the Frobenius eigenvalues $\alpha_i$. The formalism of \cite{de73} 
applies, since $\sF(\psi)$ and $\eR^{n-1} \pi_! \sF(\psi\sigma)$ 
belong to infinite compatible systems of $\ell$-adic 
representations. This reduces us to local problems, which can be 
simplified by noting that (cf. \cite[9.5]{de73}) 
\begin{enumerate}[\indent a)]
  \item the global constants for 
    $\sF(\psi)$ and $\tau^\ast \eR^{n-1} \pi_! \sF(\psi\sigma)$ are easy to 
    calculated easily, since the cohomology of these sheaves is elementary. 
  \item at every point, for either of these sheaves, the 
    corresponding $\ell$-adic representation of the decomposition group 
    has unramified semisimplification. 
\end{enumerate}

The result, for the sum with $n$ variables, is that 
\begin{equation*}\tag{7.15.2}\label{VI:eq:7-15-2}
  \det\left(F,\h_c^{n-1}(V_a^{n-1},\sF(\psi))\right) = q^{\frac{n(n-1)}{2}} \text{.} 
\end{equation*}





\subsection{Remark}\label{VI:7-16}

For $p=2$, the sheaf $\sf(\psi)$ is orthogonal. For $a\ne 0$ and $n$ 
odd, we deduce from \ref{VI:7-4}(ii) and Poincaré duality a 
symmetric bilinear form on $\h_c^{n-1}(V_a^{n-1},\sF(\psi))$ with 
values in $\bar\dQ_\ell(1-n)$. With respect to this form, $F$ is an 
orthogonal similitude of multiplier $q^{n-1}$, and $q^{\frac{1-n}{2}} F$ 
belongs to the special orthogonal group, by \eqref{VI:eq:7-15-2}. One 
of the eigenvalues of $F$ is therefore $q^{\frac{n-1}{2}}$, and the others 
occur in pairs of roots $\alpha$, $\bar\alpha$ whose product is $q^{n-1}$. 

For $n=3$, L.\ Carlitz \cite{ca69-2} obtained a more precise result, 
equivalent to the following proposition. 





\begin{proposition_}\label{VI:7-17}
For $p=2$, and $\psi(x)=(-1)^{\tr_{\dF_q/\dF_2}(x)}$, we have an isomorphism 
\[
  \h_c^2(V_a^2,\sF(\psi\sigma)) = \operatorname{Sym}^2 \h_c^1(V_a^1,\sF(\psi\sigma)) \text{.} 
\]
\end{proposition_}

The second tensor power of $\h_c^1(V_a^1,\sF(\psi\sigma))$ is 
$\h_c^2(V_a^1\times V_a^1,\sF(\psi\sigma)\boxtimes \sF(\psi\sigma))$ 
(K\"unneth), with the symmetry $x\otimes y\mapsto y\otimes x$ represented 
geometrically by $-\tau^\ast$, where $\tau$ is the automorphism 
$(x,y)\mapsto (y,x)$ of $V_a\times V_a$. If $X'$ is the quotient of 
$V_a^1\times V_a^1$ by $\{\operatorname{id}{},\tau\}$ and $\pi$ is the 
projection of $V_a^1\times V_a^1$ onto $X'$, the group 
$\operatorname{Sym}^2 \h_c^1(V_a^1,\sF(\psi\sigma))$ is therefore identified with the 
cohomology of $X'$ with coefficients in the part anti-invariant under $\tau$ of 
$\pi_\ast(\sF(\psi\sigma)\boxtimes \sF(\psi\sigma))$. Denote by $s$ the function 
on $X'$ such that $s\pi =\sigma\pr_1+\sigma\pr_2$. We have 
$\sF(\psi\sigma)\boxtimes \sF(\psi\sigma) = \sF(\psi(\sigma\pr_1+\sigma\pr_2)) = \sF(\psi s \pi) = \pi^\ast \sF(\psi s)$ 
(an isomorphism compatible with $\tau$). On the image $\pi(\Delta)$ of the 
diagonal, the desired sheaf of anti-invariants is therefore zero, while on 
the complement $X=X'\setminus \pi(\Delta)$ it is the tensor product of 
$\sF(\psi s)$ by $\varepsilon(P)$, where $\varepsilon$ is the unique character 
of order $2$ of $\{1,\tau\}$, and $P$ the $\{1,\tau\}$-torsor on $X$ given by the 
double covering induced by $V_a^1\times V_a^1$. We have 
\begin{equation*}\tag{7.17.1}\label{VI:eq:7-17-1}
  \operatorname{Sym}^2 \h_c^1(V_a^1,\sF(\psi\sigma)) = \h_c^2(X,\sF(\psi s)\otimes \varepsilon(P)) \text{.}
\end{equation*}

Let us compute $X$, $s$ and $\varepsilon(P)$. On $V_a^1\times V_a^1$, put 
$x_1=x_1\circ\pr_1$, $x_2=x_2\circ\pr_2$, $x_1'=x_1\circ \pr_2$, 
$x_2'=x_2\circ \pr_1$. We have $x_1 x_2=x_1'x_2'=a$. Identify functions on 
$X'$ with functions on $X$ invariant under $\tau$. Let $s_1=x_1+x_1'$, 
$s_2=x_2+x_2'$, $p_1=x_1 x_1'$, $p_2=x_2 x_2'$. We have $s_1 p_2 = a s_2$, 
$s_2 p_1 = a s_1$ and $X$ is defined in $X'$ by $s_1\ne 0$, $s_2\ne 0$. On 
$X$, we therefore have $p_1 = a s_1 s_2^{-1}$, and $(s_1,s_2)$ identifies $X$ with 
$\dG_m\times \dG_m$. On $X$, the covering $P$ has equation 
$T^2-s_1 T+a s_1 s_2^{-1}$ (take $T$ to be the coordinate $x_1$). If 
$T=s_1 t$, this is rewritten $t^2-t+a s_1^{-1} s_2^{-1}$. We therefore have 
$\varepsilon(P) = \sF(\psi(a s_1^{-1} s_2^{-1}))$. Since $s=s_1+s_2$, we have 
$\sF(\psi s)\otimes \varepsilon(P) = \sF(\psi(s_1+s_2+ a s_1^{-1} s_2^{-1}))$ 
and the right-hand side of \eqref{VI:eq:7-17-1} is identified with the left-hand side of 
\ref{VI:7-17}. 


\paragraph{Remark}

$\h_c^1(V_a^1,\sF(\psi\sigma))$ is also the first cohomology group of 
the elliptic curve obtained by completing the affine curve with equations 
\begin{align*}
  y^2-y &= x_1+x_2 \\
  x_1 x_2 &= a \text{;} 
\end{align*}
its modular invariant is $j=a^{-2}$. 





\subsection{Remark}\label{VI:7-18}

For $\chi$ a multiplicative character of $\dF_q^\times$, and $a\ne 0$, the 
method of \ref{VI:7-14} applies to the study of the sum 
\[
  S = \sum_{x_0\dotsm x_n = a} \chi(x_0) \psi(x_0+\cdots + x_n) 
\]
(replace the sheaf $\sF(\psi)$ on $\dG_m$ by $\sF(\chi\cdot \psi)$). We 
again find a cohomology group of dimension $n+1$, and 
$|S|\leqslant (n+1)q^{n/2}$. Such sums were considered by 
Salie and Mordell \cite{mo73}.

One may hope that our methods make it possible to study the sums 
\[
  S_{k,\chi,a} = \sum_{N(x)=a} \chi(x) \psi \tr(x) 
\]
($k$ an étale algebra over $\dF_q$, $a\in \dF_q^\times$, $\chi$ a character of 
$k^\times$). The analog of \ref{VI:7-4} (for $a\ne 0$) should be true. 





\subsection{}\label{VI:7-19}

The symmetric group $S_n$ acts on $\dA^n$, preserving $V_a$ and 
the map $\sigma$, and hence also the sheaf $\sf(\psi\sigma)$, the inverse 
image under $\sigma$ of the sheaf $\sf(\psi)$ on $\dG_a$. It therefore acts 
on $\h_c^{n-1}(V_a,\sF(\psi\sigma))$. Formula \eqref{VI:eq:7-2-5} admits 
the following cohomological interpretation.





\begin{proposition_}\label{VI:7-20}
The action of $S_n$ on $\h_c^{n-1}(V_a,\sF(\psi\sigma))$ is multiplication 
by the sign character. 
\end{proposition_}

It suffices to show that a transposition $\tau$ acts by multiplication 
by $-1$. Since $\tau^2=1$, the only possible eigenvalues of $\tau$ are 
$\pm 1$. Let $V_a/\tau$ be the quotient of $V_a$ for 
$\{\operatorname{id},\tau\}$. The map $\sigma$, and hence the sheaf 
$\sf(\psi\sigma)$, descend to the quotient, and the subspace of 
  $\h_c^{n-1}(V_a,\sF(\psi\sigma))$ fixed by 
$\tau$ is $\h_c^{n-1}(V_a/\tau,\sF(\psi\sigma))$. We must prove that 
this cohomology is zero. 

Suppose $a\ne 0$. For $\tau=(1,2)$ the quotient $V_a/\tau$ is identified 
with $\dG_a\times \dG_m^{n-2}$, the quotient map being 
$(x_1,x_2,x_3,\dots,x_n)\mapsto (x_1+x_2,x_3,\dots,x_n)$. Indeed, $x_1$ and 
$x_2$ are determined up to order by $x_1+x_2$ and 
$x_1x_2 = a/x_3\dotsm x_n$. In the coordinates $(t,x_3,\dots,x_n)$, the 
sheaf takes the form 
$\sF(\psi(t+x_3+\cdots + x_n)) = \sF(\psi(t))\otimes \sF(\psi(x_3+\cdots + x_n))$. 
  The K\"unneth formula and $\h_c^\bullet(\dG_a,\sF(\psi)) = 0$ therefore give 
the required vanishing. 

For $a=0$, we can deduce \ref{VI:7-20} from calculation \ref{VI:7-7}. 





\subsection{}\label{VI:7-21}

The sums $\kloos_{n,a}$ in this section were studied by 
S.\ Sperber, in his thesis, and in \cite{sp77} by methods 
$p$-adic methods inspired by Dwork and by Bombieri's article \cite{bo66}. 
For $p\ne 2$, these methods determine the dependence on $q$ in \ref{VI:7-5}, 
the functional equation $\alpha_i = q^{n-1} \beta_i^{-1}$ between the 
Frobenius eigenvalues for the sums $\kloos_{n,a}$ and 
$\kloos_{n,b}$ if $b=(-1)^n a$ (see \ref{VI:7-6}), and the value \ref{VI:7-15} 
of the product of the roots. For $p\geqslant n+3$, they also determine the 
$p$-adic valuations of the roots. The result is very beautiful: if we arrange 
the roots $\alpha_i$ in a suitable order, with 
  $0\leqslant i<n$, $\alpha_i q^{-i}$ are units. The 
restrictions on $p$ are probably not essential. 










\section{Other applications}\label{VI:8}





\subsection{}\label{VI:8-1}

The identity \ref{VI:1-9-4} can sometimes be used to calculate 
the number of rational points of an algebraic variety over $\dF_q$. In 
most cases where an exact result has been obtained, cohomology 
is expressed in terms of algebraic cycles; this is the case for rational 
surfaces, quadrics, smooth intersections of two quadrics of 
odd dimension, flag varieties\ldots. This is not, however, 
always the case; in Lusztig's beautiful article \cite{lu76}, algebraic cycles 
do not appear. 





\subsection{}\label{VI:8-2}

For reductive groups, cohomology can be calculated by lifting to 
characteristic $0$: if $B$ is a Borel subgroup of $G$, $G$ is a 
$B$-torsor over $G/B$; the cohomology of the smooth projective variety 
$G/B$ is invariant under specialization, and we use the Leray spectral sequence 
of $G\to G/B$ to prove the same result for $G$. It can be expressed 
conveniently in terms of \emph{the} maximal torus $T$ of $G$ and the action 
\emph{of the} Weyl group $W$ on $T$ (for the definition \emph{of the} maximal 
torus, cf. \cite[VIII \S 2 Rem.2]{bo05} or \cite[p.105]{dl76}). It is 
the exterior algebra on its primitive part, which agrees up to a shift 
with the quotient $I(\h^\bullet(B G,\dQ_\ell))$ of the cohomology of 
$B G$ formed by the ``indecomposable elements,'' and 
$\h^\bullet(B G,\dQ_\ell)=\h^\bullet(B T,\dQ_\ell)^W$. If $X$ is the group of 
characters of $T$, we have altogether 
\[
  \h^\bullet(G,\dQ_\ell) = \bigwedge \left(I\left(\operatorname{Sym}^\bullet(X\otimes \dQ_\ell(-1))^W\right)[-1]\right) \text{.}
\]
If $G$ is defined over a finite field $\dF_q$, $\gal(\dF/\dF_q)$ acts on $X$ 
and this isomorphism is Galois-compatible. Passing to 
cohomology with compact supports, we recover the classical formulas for the 
number of rational points of $G$. If $F^\ast:X\to X$ is defined by the 
Frobenius morphism $F:T\to T$, and $\deg_G(F)$ is the degree 
$q^{\dim G}$ of $F:G\to G$, we have 
\begin{equation*}\tag{8.2.1}\label{VI:eq:8-2-1}
  \left|G_0(\dF_q)\right| = \deg_G(F) \cdot \det\left(1-{F^\ast}^{-1},I(\operatorname{Sym}^\bullet(X\otimes \dQ)^W)\right) \text{.}
\end{equation*}

The same formula applies to the Ree and Suzuki groups; these groups are 
of the form $G^F$ (reductive $G$, $F^2$ a Frobenius), and it suffices to 
verify that the Lefschetz trace formula holds for $F$ acting 
on $G$. If $B$ is a Borel subgroup stable by $F$, we verify it 
directly for $F$ acting on $B$; for the action on $G/B$, which is proper and 
smooth, we can invoke the general theorems. 





\subsection{}\label{VI:8-3}

The trace formula was used by Deligne--Lusztig \cite{dl76}, 
Kazhdan \cite{ka77} and Springer \cite{sp76} in the study of 
complex representations of finite groups $G_0(\dF_q)$, for $G_0$ 
reductive over $\dF_q$. The work of Kazhdan and Springer contains 
admirable examples of how the trace formula permits ``analytic 
continuation'' from a split situation to a nonsplit situation 
(see \ref{VI:1-13}).
